\chapter{Coequational Subobjects and Birkhoff Theorems}
\label{ch:coequational-subobjects-birkhoff}

\section{Overview and standing assumptions}
\label{sec:coeq-birkhoff-overview}

\subsection{Orientation}
\label{subsec:coeq-birkhoff-orientation}

The preceding chapter constructed, for a behaviour specification
\[
    k:K\to \mathbf{Beh}_0,
\]
its derivative image
\[
    \operatorname{Der}_0(K)\hookrightarrow \mathbf{Beh}_0
\]
and the corresponding minimal residual Mealy morphism
\[
    \mathsf{Min}(K):\mathbb A\rightsquigarrow\mathbb B.
\]
The present chapter reorganizes residual semantics in coequational form. The
central semantic object is the canonical behaviour Mealy morphism
\[
    \mathbf{Beh}^{\mathrm{can}}_{\mathbb A,\mathbb B}:
    \mathbb A\rightsquigarrow\mathbb B.
\]
By Chapter~\ref{ch:behaviour-categories-myhill-nerode}, this object is
terminal in the ordinary category
\[
    \mathsf{Mly}^{\mathrm{str}}_{\mathbb A,\mathbb B}
\]
of Mealy morphisms and strict semantic homomorphisms. Hence a local
coequational theory is a subobject of this terminal semantic object:
\[
    V\hookrightarrow \mathbf{Beh}^{\mathrm{can}}_{\mathbb A,\mathbb B}.
\]
A Mealy morphism
\[
    P:\mathbb A\rightsquigarrow\mathbb B
\]
models \(V\) when its unique semantic map to the terminal object factors
through \(V\):
\[
\begin{tikzcd}[column sep=large, row sep=large]
    P
        \arrow[rr, "\beta_P"]
        \arrow[dr, dashed, "\bar\beta_P"']
    &&
    \mathbf{Beh}^{\mathrm{can}}_{\mathbb A,\mathbb B}
    \\
    &
    V
        \arrow[ur, hook, "j_V"']
    &
\end{tikzcd}
\]
Equivalently, the carrier map
\[
    \beta_P:P\to\mathbf{Beh}_0
\]
factors through the underlying carrier subobject
\[
    V\hookrightarrow\mathbf{Beh}_0
\]
in the slice over
\[
    Z:=A_0\times B_0.
\]

The basic examples of coequational theories are semantic images. For a Mealy
morphism \(P\), its semantic image is the regular image of its terminal
semantic map:
\[
    P\twoheadrightarrow \operatorname{Beh}(P)
    \hookrightarrow \mathbf{Beh}^{\mathrm{can}}_{\mathbb A,\mathbb B}.
\]
The first main task of the chapter is to prove that regular images of strict
semantic homomorphisms are created by the ambient carrier slice. This makes
\[
    \operatorname{Beh}(P)
\]
again a subobject of the canonical behaviour Mealy morphism.

The chapter proves a coequational Birkhoff theorem. Locally, a replete full
class of Mealy morphisms is definable by a local coequational theory exactly
when it is closed under strict homomorphic preimages, regular-epimorphic
quotients, and small coproducts. The representing coequational theory
associated to a class
\[
    \mathcal C\subseteq\mathsf{Mly}^{\mathrm{str}}_{\mathbb A,\mathbb B}
\]
is the join of the semantic images of its objects:
\[
    V_{\mathcal C}
    :=
    \bigvee_{P\in\mathcal C}\operatorname{Beh}(P)
    \hookrightarrow \mathbf{Beh}_0.
\]
The use of arbitrary joins of subobjects is the point at which this chapter
works in a Grothendieck topos. The resulting statement is in the same lineage
as Birkhoff and Eilenberg type correspondences for recognizable languages and
categorical varieties of languages
\cite{Eilenberg1974,AdamekMiliusMyersUrbat2019TOCL,Salamanca2016Monad,Salamanca2017Unveiling}.

\subsection{Ambient assumptions}
\label{subsec:coeq-birkhoff-ambient-assumptions}

\begin{assumption}
\label{ass:coeq-birkhoff-grothendieck-topos}
    Throughout this chapter, \(\mathcal E\) is a Grothendieck topos with chosen
    pullbacks. Thus \(\mathcal E\) is Barr-exact and locally cartesian closed,
    has stable regular epi--mono factorizations, effective quotients of
    internal equivalence relations, small coproducts, and complete subobject
    lattices in every slice. All coproducts and joins below are small, relative
    to the fixed ambient universe. We use these topos-theoretic properties in
    the form recorded in
    \cite{Johnstone2002,Borceux1994,Jacobs1999CategoricalLogic}.
\end{assumption}

The input and output internal categories are
\[
    \mathbb A=(A_0,A_1,s_A,t_A,e_A,m_A),
    \qquad
    \mathbb B=(B_0,B_1,s_B,t_B,e_B,m_B).
\]
When the input and output categories are fixed, we write
\[
    Z:=A_0\times B_0.
\]
Composition in \(\mathbb A\) is written
\[
    m_A(a,b)=b\circ a
\]
for composable arrows
\[
    a:u\to v,
    \qquad
    b:v\to w.
\]
Thus a Mealy state over \(w\) processes \(b\circ a\) by first processing \(b\),
then processing \(a\).

Small coproducts include the empty coproduct. Thus every coequationally
definable class contains the initial strict Mealy morphism. Dually, the empty
join of local coequational theories is the bottom subobject of the canonical
behaviour object.

\subsection{Strict Mealy morphisms and terminal semantics}
\label{subsec:coeq-standing-notation-terminal-semantics}

Throughout, Mealy morphisms
\[
    P:\mathbb A\rightsquigarrow\mathbb B
\]
are considered as objects of the ordinary category
\[
    \mathsf{Mly}^{\mathrm{str}}_{\mathbb A,\mathbb B}.
\]
A strict semantic homomorphism
\[
    f:P\to Q
\]
is a map in the carrier slice
\[
    \mathcal E/Z
\]
preserving transition and output on the nose. Thus, writing
\[
    P=(P,p_P,q_P,\delta_P,\omega_P),
    \qquad
    Q=(Q,p_Q,q_Q,\delta_Q,\omega_Q),
\]
strictness means
\[
    p_Qf=p_P,
    \qquad
    q_Qf=q_P,
\]
and
\[
    f\delta_P
    =
    \delta_Q(1\times f),
    \qquad
    \omega_P
    =
    \omega_Q(1\times f).
\]
Subobjects, regular quotients, coproducts, and full subcategories below are
therefore taken in
\[
    \mathsf{Mly}^{\mathrm{str}}_{\mathbb A,\mathbb B}
\]
unless explicitly stated otherwise. A regular quotient means a
regular-epimorphic strict semantic homomorphism.

The strictness equations may be displayed by
\[
\begin{tikzcd}[column sep=huge, row sep=large]
    A_1\times_{A_0}P
        \arrow[r, "{1\times f}"]
        \arrow[d, "\delta_P"']
    &
    A_1\times_{A_0}Q
        \arrow[d, "\delta_Q"]
    \\
    P
        \arrow[r, "f"']
    &
    Q
\end{tikzcd}
\qquad
\begin{tikzcd}[column sep=huge, row sep=large]
    A_1\times_{A_0}P
        \arrow[r, "{1\times f}"]
        \arrow[dr, "\omega_P"']
    &
    A_1\times_{A_0}Q
        \arrow[d, "\omega_Q"]
    \\
    &
    B_1 .
\end{tikzcd}
\]
The carrier condition is
\[
\begin{tikzcd}[column sep=large, row sep=large]
    P
        \arrow[rr, "f"]
        \arrow[dr, "{(p_P,q_P)}"']
    &&
    Q
        \arrow[dl, "{(p_Q,q_Q)}"]
    \\
    &
    Z .
\end{tikzcd}
\]

Recall from Chapter~\ref{ch:behaviour-categories-myhill-nerode} that every
Mealy morphism \(P\) has a residual semantics map
\[
    \beta_P:P\to\mathbf{Beh}_0.
\]
For \(x\in P\), the residual behaviour
\[
    \beta_P(x)=H_x:\mathbb A/p_P(x)\to\mathbb B
\]
is given on objects by
\[
    H_x(r)=q_P(\delta_P(r,x)),
\]
and on residual arrows
\[
    w\xrightarrow{a}u\xrightarrow{r}p_P(x)
\]
by
\[
    H_x(a,r)=\omega_P(a,\delta_P(r,x)).
\]
Equivalently,
\[
    \beta_P:P\to\mathbf{Beh}^{\mathrm{can}}_{\mathbb A,\mathbb B}
\]
is the unique strict semantic homomorphism from \(P\) to the terminal object
\[
    \mathbf{Beh}^{\mathrm{can}}_{\mathbb A,\mathbb B}
\]
of \(\mathsf{Mly}^{\mathrm{str}}_{\mathbb A,\mathbb B}\). The use of terminal
semantic objects follows the categorical semantics of realization and
coalgebraic behaviour maps
\cite{Goguen1973,Rutten2000,Jacobs2016Coalgebra}.

\begin{lemma}[Strict maps preserve residual semantics]
\label{lem:coeq-strict-maps-preserve-beta}
    Let
    \[
        f:P\to Q
    \]
    be a strict semantic homomorphism. Then
    \[
        \beta_Qf=\beta_P.
    \]
\end{lemma}

\begin{proof}
    \leavevmode
    \begin{enumerate}
        \item \textbf{Use terminality.}
        Since \(f\) is strict and \(\beta_Q\) is strict, the composite
        \[
            \beta_Qf:P\to\mathbf{Beh}^{\mathrm{can}}_{\mathbb A,\mathbb B}
        \]
        is a strict semantic homomorphism. The map \(\beta_P\) is also a strict
        semantic homomorphism from \(P\) to the terminal object.

        This gives the parallel pair
        \[
        \begin{tikzcd}[column sep=huge, row sep=large]
            P
                \arrow[r, shift left=0.65ex, "\beta_Qf"]
                \arrow[r, shift right=0.65ex, "\beta_P"']
            &
            \mathbf{Beh}^{\mathrm{can}}_{\mathbb A,\mathbb B}.
        \end{tikzcd}
        \]

        \item \textbf{Conclude uniqueness.}
        By terminality of
        \[
            \mathbf{Beh}^{\mathrm{can}}_{\mathbb A,\mathbb B}
        \]
        in \(\mathsf{Mly}^{\mathrm{str}}_{\mathbb A,\mathbb B}\), the two maps
        agree:
        \[
            \beta_Qf=\beta_P.
        \]
    \end{enumerate}
\end{proof}

\subsection{Running specializations}
\label{subsec:coeq-overview-running-specializations}

The two running specializations from the previous chapters will continue to be
used.

First, in the one-object case, the input internal category is determined by a
monoid object \(M\), and the output internal category is determined by a monoid
object \(N\), with the convention
\[
    b\circ a=ba.
\]
A Mealy morphism is an internal right \(M\)-object together with an
\(N\)-valued output cocycle. In the Boolean language specialization, where
\[
    \mathcal E=\mathbf{Set},
    \qquad
    M=\Sigma^*,
    \qquad
    \mathbb B=2_{\mathrm{ch}},
\]
a residual behaviour is a language
\[
    L:\Sigma^*\to 2,
\]
and the derivative convention is
\[
    (\partial_uL)(w)=L(uw).
\]

Second, in the stateless case, a Mealy morphism is induced by an internal
functor
\[
    F:\mathbb A\to\mathbb B.
\]
Its residual behaviour at \(v\in A_0\) is the functor
\[
    H^F_v:\mathbb A/v\to\mathbb B
\]
given by
\[
    H^F_v(r:u\to v)=F_0(u),
    \qquad
    H^F_v(a,r)=F_1(a).
\]
Thus a coequational theory contains the stateless recognizer \(F\) exactly when
all residual functors
\[
    H^F_v
\]
belong to the chosen behaviour subobject.

\section{Local coequational theories}
\label{sec:local-coequational-theories}

\subsection{Canonical derivative structure}
\label{subsec:canonical-derivative-structure}

The canonical behaviour object
\[
    \mathbf{Beh}_0
\]
carries the canonical behaviour Mealy morphism
\[
    \mathbf{Beh}^{\mathrm{can}}_{\mathbb A,\mathbb B}:
    \mathbb A\rightsquigarrow\mathbb B.
\]
A state over \(v\in A_0\) is a residual behaviour
\[
    H:\mathbb A/v\to\mathbb B.
\]
For an arrow
\[
    a:u\to v,
\]
the canonical transition is the derivative
\[
    \partial_aH=H\circ a_*,
\]
and the canonical output comparison is
\[
    \omega_{\mathbf{Beh}}(a,H):
    q_{\mathbf{Beh}}(\partial_aH)\to q_{\mathbf{Beh}}(H),
\]
the arrow assigned by \(H\) to the residual morphism
\[
    a\to 1_v
\]
in \(\mathbb A/v\). This derivative notation is compatible with the
Brzozowski derivative viewpoint in the one-object free-monoid specialization
\cite{Brzozowski1964}.

The output comparison is displayed by applying \(H\) to the residual triangle
\[
\begin{tikzcd}[column sep=large, row sep=large]
    u
        \arrow[r, "a"]
        \arrow[rr, "a", bend left=15]
    &
    v
        \arrow[r, "1_v"']
    &
    v .
\end{tikzcd}
\qquad
\leadsto
\qquad
\begin{tikzcd}[column sep=large, row sep=large]
    H(a:u\to v)
        \arrow[r, "{H(a\to 1_v)}"]
    &
    H(1_v).
\end{tikzcd}
\]
Equivalently, the transition-output structure is the map of spans
\[
\begin{tikzcd}[column sep=large, row sep=large]
    A_1\times_{t_A,A_0,p_{\mathbf{Beh}}}\mathbf{Beh}_0
        \arrow[rr, "{\langle \delta_{\mathbf{Beh}},\omega_{\mathbf{Beh}}\rangle}"]
        \arrow[dr, "{s_A\pi_{A_1}}"']
    &&
    \mathbf{Beh}_0\times_{q_{\mathbf{Beh}},B_0,s_B}B_1
        \arrow[dl, "{p_{\mathbf{Beh}}\pi_{\mathbf{Beh}}}"]
        \arrow[dr, "{t_B\pi_{B_1}}"]
    \\
    &
    A_0
    &&
    B_0 .
\end{tikzcd}
\]

Whenever a derivative domain occurs as a carrier over
\[
    Z=A_0\times B_0,
\]
it is regarded over \(Z\) by the derivative state. Thus the object
\[
    A_1\times_{t_A,A_0,p_{\mathbf{Beh}}}\mathbf{Beh}_0
\]
is regarded as an object of \(\mathcal E/Z\) by the composite
\[
\begin{tikzcd}[column sep=large, row sep=large]
    A_1\times_{A_0}\mathbf{Beh}_0
        \arrow[r, "\delta_{\mathbf{Beh}}"]
        \arrow[dr, dashed, "{(s_A\pi_{A_1},\,q_{\mathbf{Beh}}\delta_{\mathbf{Beh}})}"']
    &
    \mathbf{Beh}_0
        \arrow[d, "{(p_{\mathbf{Beh}},q_{\mathbf{Beh}})}"]
    \\
    &
    Z .
\end{tikzcd}
\]
Thus, if
\[
    a:u\to v
\]
and \(H\) lies over \(v\), then \((a,H)\) has \(Z\)-coordinate
\[
    \bigl(u,\;q_{\mathbf{Beh}}(\partial_aH)\bigr).
\]
This convention is used in all regular image and factorization statements
involving derivative domains.

\subsection{Coequational subobjects}
\label{subsec:coequational-subobjects}

\begin{definition}[Local coequational theory]
\label{def:local-coequational-theory}
    A \textbf{local coequational theory} over \((\mathbb A,\mathbb B)\) is a
    subobject
    \[
        V\hookrightarrow \mathbf{Beh}^{\mathrm{can}}_{\mathbb A,\mathbb B}
    \]
    in
    \[
        \mathsf{Mly}^{\mathrm{str}}_{\mathbb A,\mathbb B}.
    \]
\end{definition}

The following proposition identifies this definition with the expected
derivative-closed carrier condition.

\begin{proposition}[Subobjects of the canonical behaviour machine]
\label{prop:subobjects-canonical-behaviour-machine}
    To give a subobject
    \[
        V\hookrightarrow \mathbf{Beh}^{\mathrm{can}}_{\mathbb A,\mathbb B}
    \]
    in
    \[
        \mathsf{Mly}^{\mathrm{str}}_{\mathbb A,\mathbb B}
    \]
    is equivalently to give a subobject
    \[
        j_V:V\hookrightarrow\mathbf{Beh}_0
    \]
    in \(\mathcal E/Z\) such that the canonical derivative map restricts to a
    transition
    \[
        \delta_V:
        A_1\times_{t_A,A_0,p_V}V\to V.
    \]
    Under this condition, the canonical output map restricts to
    \[
        \omega_V:
        A_1\times_{t_A,A_0,p_V}V\to B_1.
    \]
\end{proposition}

\begin{proof}
    \leavevmode
    \begin{enumerate}
        \item \textbf{Forget a strict subobject.}
        A subobject in
        \[
            \mathsf{Mly}^{\mathrm{str}}_{\mathbb A,\mathbb B}
        \]
        has an underlying monomorphism
        \[
            j_V:V\hookrightarrow\mathbf{Beh}_0
        \]
        in the carrier slice \(\mathcal E/Z\). Since the inclusion is strict,
        the transition of \(V\) is carried by the canonical transition:
        \[
            j_V\delta_V
            =
            \delta_{\mathbf{Beh}}(1\times j_V).
        \]
        Hence the canonical derivative map lands in \(V\). The displayed
        restriction is
        \[
        \begin{tikzcd}[column sep=huge, row sep=large]
            A_1\times_{A_0}V
                \arrow[r, "{1\times j_V}"]
                \arrow[d, "\delta_V"']
            &
            A_1\times_{A_0}\mathbf{Beh}_0
                \arrow[d, "\delta_{\mathbf{Beh}}"]
            \\
            V
                \arrow[r, hook, "j_V"']
            &
            \mathbf{Beh}_0 .
        \end{tikzcd}
        \]

        \item \textbf{Restrict the canonical structure.}
        Conversely, suppose
        \[
            j_V:V\hookrightarrow\mathbf{Beh}_0
        \]
        is a carrier subobject closed under the canonical derivative map. Define
        \[
            \delta_V(a,H)=\partial_aH.
        \]
        Define
        \[
            \omega_V(a,H)=\omega_{\mathbf{Beh}}(a,H).
        \]
        The typing equations are inherited from the canonical behaviour
        machine. More explicitly, the following diagram commutes:
        \[
        \begin{tikzcd}[column sep=huge, row sep=large]
            A_1\times_{A_0}V
                \arrow[r, "{1\times j_V}"]
                \arrow[d, "\delta_V"']
                \arrow[dr, phantom, "\circlearrowleft", very near start]
            &
            A_1\times_{A_0}\mathbf{Beh}_0
                \arrow[r, "\omega_{\mathbf{Beh}}"]
                \arrow[d, "\delta_{\mathbf{Beh}}"]
            &
            B_1
                \arrow[d, "{(s_B,t_B)}"]
            \\
            V
                \arrow[r, hook, "j_V"']
            &
            \mathbf{Beh}_0
                \arrow[r, "{(q_{\mathbf{Beh}},q_{\mathbf{Beh}})}"']
            &
            B_0\times B_0 ,
        \end{tikzcd}
        \]
        where the source component is read after \(\delta_{\mathbf{Beh}}\) and
        the target component is read at the original state.

        \item \textbf{Verify the Mealy laws.}
        The unit and multiplication laws for \(\delta_V\) and \(\omega_V\)
        hold after composing with the monomorphism \(j_V\), where they become
        the laws of
        \[
            \mathbf{Beh}^{\mathrm{can}}_{\mathbb A,\mathbb B}.
        \]
        Hence they hold in \(V\).

        \item \textbf{Identify the strict inclusion.}
        By construction,
        \[
            j_V:V\to\mathbf{Beh}^{\mathrm{can}}_{\mathbb A,\mathbb B}
        \]
        preserves transition and output on the nose. Therefore it is a strict
        semantic homomorphism and exhibits \(V\) as a subobject of the
        canonical behaviour machine.
    \end{enumerate}
\end{proof}

The derivative-closedness condition is the factorization
\[
\begin{tikzcd}[column sep=large, row sep=large]
    A_1\times_{t_A,A_0,p_V}V
        \arrow[rr, "{\partial(1\times j_V)}"]
        \arrow[dr, "\delta_V"']
    &&
    \mathbf{Beh}_0
    \\
    &
    V
        \arrow[ur, hook, "j_V"']
    &
\end{tikzcd}
\]
in \(\mathcal E/Z\). Equivalently, the transition-output cell of the canonical
machine restricts along \(j_V\):
\[
\begin{tikzcd}[column sep=huge, row sep=large]
    A_1\times_{A_0}V
        \arrow[r, "{1\times j_V}"]
        \arrow[d, "{\langle\delta_V,\omega_V\rangle}"']
    &
    A_1\times_{A_0}\mathbf{Beh}_0
        \arrow[d, "{\langle\delta_{\mathbf{Beh}},\omega_{\mathbf{Beh}}\rangle}"]
    \\
    V\times_{B_0}B_1
        \arrow[r, hook, "{j_V\times 1}"']
    &
    \mathbf{Beh}_0\times_{B_0}B_1 .
\end{tikzcd}
\]

When \(V\hookrightarrow\mathbf{Beh}_0\) is a local coequational theory, we write
\[
    \delta_V(a,H)=\partial_aH
\]
and
\[
    \omega_V(a,H)=\omega_{\mathbf{Beh}}(a,H).
\]

\subsection{Models of local coequational theories}
\label{subsec:models-local-coequational-theories}

\begin{definition}[Models]
\label{def:models-coequational-theory}
    Let
    \[
        V\hookrightarrow\mathbf{Beh}^{\mathrm{can}}_{\mathbb A,\mathbb B}
    \]
    be a local coequational theory. A Mealy morphism
    \[
        P:\mathbb A\rightsquigarrow\mathbb B
    \]
    \textbf{models} \(V\) if its terminal semantic map
    \[
        \beta_P:P\to\mathbf{Beh}^{\mathrm{can}}_{\mathbb A,\mathbb B}
    \]
    factors through \(V\). We write
    \[
        \mathsf{Mod}(V)
    \]
    for the full subcategory of
    \[
        \mathsf{Mly}^{\mathrm{str}}_{\mathbb A,\mathbb B}
    \]
    on the objects that model \(V\).
\end{definition}

\begin{lemma}[Model factorizations are strict]
\label{lem:model-factorizations-strict}
    If \(P\models V\), then the unique factorization
    \[
        \bar\beta_P:P\to V
    \]
    satisfying
    \[
        j_V\bar\beta_P=\beta_P
    \]
    is a strict semantic homomorphism.
\end{lemma}

\begin{proof}
    \leavevmode
    \begin{enumerate}
        \item \textbf{Check the map over \(Z\).}
        The equation
        \[
            j_V\bar\beta_P=\beta_P
        \]
        and the fact that both \(j_V\) and \(\beta_P\) are maps over \(Z\)
        imply that \(\bar\beta_P\) is a map over \(Z\). This is the lower
        triangle in
        \[
        \begin{tikzcd}[column sep=large, row sep=large]
            P
                \arrow[rr, "\beta_P"]
                \arrow[dr, "\bar\beta_P"']
                \arrow[ddr, "{(p_P,q_P)}"', bend right=18]
            &&
            \mathbf{Beh}_0
                \arrow[ddl, "{(p_{\mathbf{Beh}},q_{\mathbf{Beh}})}", bend left=18]
            \\
            &
            V
                \arrow[ur, hook, "j_V"']
                \arrow[d, "{(p_V,q_V)}"]
            &
            \\
            &
            Z .
        \end{tikzcd}
        \]

        \item \textbf{Check transition preservation.}
        Both composites
        \[
            j_V\bar\beta_P\delta_P
            \qquad\text{and}\qquad
            j_V\delta_V(1\times\bar\beta_P)
        \]
        are equal to the transition-preservation composite for the strict map
        \[
            \beta_P:P\to\mathbf{Beh}^{\mathrm{can}}_{\mathbb A,\mathbb B}.
        \]
        The comparison is the outer rectangle
        \[
        \begin{tikzcd}[column sep=huge, row sep=large]
            A_1\times_{A_0}P
                \arrow[r, "{1\times \bar\beta_P}"]
                \arrow[d, "\delta_P"']
                \arrow[rr, bend left=16, "{1\times \beta_P}"]
            &
            A_1\times_{A_0}V
                \arrow[r, "{1\times j_V}"]
                \arrow[d, "\delta_V"]
            &
            A_1\times_{A_0}\mathbf{Beh}_0
                \arrow[d, "\delta_{\mathbf{Beh}}"]
            \\
            P
                \arrow[r, "\bar\beta_P"']
                \arrow[rr, bend right=16, "\beta_P"']
            &
            V
                \arrow[r, hook, "j_V"']
            &
            \mathbf{Beh}_0 .
        \end{tikzcd}
        \]
        Since \(j_V\) is monic, the transition equation holds in \(V\).

        \item \textbf{Check output preservation.}
        The output equation for \(\bar\beta_P\) follows from the output equation
        for \(\beta_P\), since \(\omega_V\) is the restriction of
        \(\omega_{\mathbf{Beh}}\). It is represented by
        \[
        \begin{tikzcd}[column sep=huge, row sep=large]
            A_1\times_{A_0}P
                \arrow[r, "{1\times \bar\beta_P}"]
                \arrow[dr, "\omega_P"']
            &
            A_1\times_{A_0}V
                \arrow[d, "\omega_V"]
                \arrow[r, "{1\times j_V}"]
            &
            A_1\times_{A_0}\mathbf{Beh}_0
                \arrow[dl, "\omega_{\mathbf{Beh}}"]
            \\
            &
            B_1 .
        \end{tikzcd}
        \]
    \end{enumerate}
\end{proof}

\subsection{Restricted canonical subobjects}
\label{subsec:restricted-canonical-subobjects}

\begin{lemma}[Semantic map of a restricted canonical subobject]
\label{lem:semantic-map-of-a-restricted-canonical-subobject}
    Let
    \[
        V\hookrightarrow\mathbf{Beh}^{\mathrm{can}}_{\mathbb A,\mathbb B}
    \]
    be a local coequational theory. Then the residual semantics map
    \[
        \beta_V:V\to\mathbf{Beh}_0
    \]
    is the inclusion
    \[
        j_V:V\hookrightarrow\mathbf{Beh}_0.
    \]
\end{lemma}

\begin{proof}
    \leavevmode
    \begin{enumerate}
        \item \textbf{Compute residual objects.}
        Let
        \[
            H\in V_v
        \]
        and let
        \[
            r:u\to v
        \]
        be a residual object. In the restricted canonical object, execution of
        \(r\) from \(H\) gives
        \[
            \partial_rH.
        \]
        Its anchor output is
        \[
            q_V(\partial_rH)=H(r).
        \]
        This is the object part of the equality displayed by
        \[
        \begin{tikzcd}[column sep=large, row sep=large]
            H
                \arrow[r, maps to]
                \arrow[d, "{\partial_r}"']
            &
            H(1_v)
            \\
            \partial_rH
                \arrow[r, maps to]
            &
            H(r).
        \end{tikzcd}
        \]

        \item \textbf{Compute residual arrows.}
        For a residual arrow
        \[
            w\xrightarrow{a}u\xrightarrow{r}v,
        \]
        the output produced by the restricted canonical object at
        \(\partial_rH\) along \(a\) is
        \[
            \omega_V(a,\partial_rH).
        \]
        This is the arrow assigned by \(H\) to the residual morphism
        \[
            r\circ a\to r.
        \]
        The comparison is
        \[
        \begin{tikzcd}[column sep=large, row sep=large]
            w
                \arrow[r, "a"]
                \arrow[rr, "r\circ a", bend left=15]
            &
            u
                \arrow[r, "r"']
            &
            v
        \end{tikzcd}
        \qquad
        \leadsto
        \qquad
        \begin{tikzcd}[column sep=large, row sep=large]
            H(r\circ a)
                \arrow[r, "{H(r\circ a\to r)}"]
            &
            H(r).
        \end{tikzcd}
        \]

        \item \textbf{Conclude equality of represented behaviours.}
        Thus the residual behaviour computed from \(H\) in the restricted
        canonical object agrees with \(H\) on residual objects and residual
        arrows. Hence
        \[
            \beta_V(H)=H.
        \]
        Therefore
        \[
            \beta_V=j_V.
        \]
        Equivalently, the terminal semantic triangle is
        \[
        \begin{tikzcd}[column sep=large, row sep=large]
            V
                \arrow[rr, "\beta_V"]
                \arrow[dr, equal]
            &&
            \mathbf{Beh}^{\mathrm{can}}
            \\
            &
            V
                \arrow[ur, hook, "j_V"']
            &
        \end{tikzcd}
        \]
        and the diagonal identity is forced by the previous computation.
    \end{enumerate}
\end{proof}

\subsection{Running specializations}
\label{subsec:local-coeq-running-specializations}

In the one-object free-monoid Boolean specialization, a local coequational
theory is a collection
\[
    V\subseteq 2^{\Sigma^*}
\]
of languages closed under residual derivatives:
\[
    L\in V
    \quad\Longrightarrow\quad
    \partial_uL\in V
    \qquad
    (u\in\Sigma^*).
\]
The model condition says that every residual language generated by a state of
the machine lies in \(V\). Thus, for a deterministic Boolean Mealy recognizer
with state set \(P\), the condition is
\[
    \{\beta_P(x)\mid x\in P\}\subseteq V.
\]

In the stateless specialization, an internal functor
\[
    F:\mathbb A\to\mathbb B
\]
models \(V\) precisely when
\[
    H^F_v\in V
\]
for every object \(v\in A_0\), where
\[
    H^F_v(r:u\to v)=F_0(u),
    \qquad
    H^F_v(a,r)=F_1(a).
\]
Thus local coequations constrain the residual slice functors generated by
\(F\), rather than merely the object map or arrow map of \(F\) separately.

\section{Semantic images and exact structure}
\label{sec:regular-semantic-images}

\subsection{Regular-epimorphic descent for subobjects}
\label{subsec:regular-epi-descent-subobjects}

The following elementary descent fact will be used repeatedly.

\begin{lemma}[Regular-epimorphic descent for subobject membership]
\label{lem:regular-epi-descent-for-subobjects-membership}
    Let \(\mathcal X\) be either \(\mathcal E\) or a slice
    \(\mathcal E/S\). Let
    \[
        e:X\twoheadrightarrow Y
    \]
    be a regular epimorphism in \(\mathcal X\), and let
    \[
        j:V\hookrightarrow W
    \]
    be a monomorphism in \(\mathcal X\). If
    \[
        h:Y\to W
    \]
    is a morphism in \(\mathcal X\) such that
    \[
        he
    \]
    factors through \(j\), then \(h\) factors uniquely through \(j\).
\end{lemma}

\begin{proof}
    \leavevmode
    \begin{enumerate}
        \item \textbf{Use the kernel pair of the regular epimorphism.}
        Let
        \[
            R\rightrightarrows X
        \]
        be the kernel pair of \(e\) in \(\mathcal X\). Since \(e\) is a
        regular epimorphism, it is the coequalizer of this kernel pair:
        \[
        \begin{tikzcd}[column sep=large, row sep=large]
            R
                \arrow[r, shift left=0.65ex, "r_1"]
                \arrow[r, shift right=0.65ex, "r_2"']
            &
            X
                \arrow[r, two heads, "e"]
            &
            Y .
        \end{tikzcd}
        \]

        \item \textbf{Descend the factorization.}
        Suppose
        \[
            he=jg
        \]
        for some \(g:X\to V\). Since
        \[
            he r_1=he r_2,
        \]
        we have
        \[
            jg r_1=jg r_2.
        \]
        Since \(j\) is monic,
        \[
            g r_1=g r_2.
        \]
        Therefore \(g\) descends uniquely through \(e\) to a map
        \[
            \bar g:Y\to V
        \]
        with
        \[
            \bar g e=g.
        \]
        The descent diagram is
        \[
        \begin{tikzcd}[column sep=large, row sep=large]
            R
                \arrow[r, shift left=0.65ex, "r_1"]
                \arrow[r, shift right=0.65ex, "r_2"']
            &
            X
                \arrow[r, two heads, "e"]
                \arrow[dr, "g"']
            &
            Y
                \arrow[d, dashed, "\bar g"]
            \\
            &&
            V .
        \end{tikzcd}
        \]

        \item \textbf{Compare with \(h\).}
        Now
        \[
            j\bar g e=jg=he.
        \]
        Since \(e\) is epic,
        \[
            j\bar g=h.
        \]
        Uniqueness follows from monicity of \(j\). The final comparison is
        \[
        \begin{tikzcd}[column sep=large, row sep=large]
            X
                \arrow[r, two heads, "e"]
                \arrow[dr, "g"']
                \arrow[ddr, "he"', bend right=18]
            &
            Y
                \arrow[d, dashed, "\bar g"]
                \arrow[dd, "h", bend left=18]
            \\
            &
            V
                \arrow[d, hook, "j"]
            \\
            &
            W .
        \end{tikzcd}
        \]
    \end{enumerate}
\end{proof}

\subsection{Regular images of strict homomorphisms}
\label{subsec:regular-images-strict-homomorphisms}

This section proves that regular images of strict semantic homomorphisms are
created by the carrier slice.

\begin{proposition}[Regular images are created by the carrier slice]
\label{prop:regular-image-factorization-mealy}
    Let
    \[
        f:P\to Q
    \]
    be a strict semantic homomorphism. Factor the underlying map in
    \(\mathcal E/Z\) as
    \[
        P\xrightarrow{e}I\xrightarrow{m}Q,
    \]
    where \(e\) is a regular epimorphism and \(m\) is a monomorphism. Then \(I\)
    carries a unique Mealy structure such that
    \[
        e:P\twoheadrightarrow I
    \]
    and
    \[
        m:I\hookrightarrow Q
    \]
    are strict semantic homomorphisms.
\end{proposition}

\begin{proof}
    \leavevmode
    \begin{enumerate}
        \item \textbf{Pull back the image cover to transition domains.}
        Since regular epimorphisms are stable under pullback, the induced map
        \[
            \widehat e:
            A_1\times_{t_A,A_0,p_P}P
            \twoheadrightarrow
            A_1\times_{t_A,A_0,p_I}I
        \]
        is a regular epimorphism. It sends
        \[
            (a,x)\longmapsto(a,ex).
        \]
        It is the vertical arrow in the pullback square
        \[
        \begin{tikzcd}[column sep=large, row sep=large]
            A_1\times_{A_0}P
                \arrow[r, "\pi_P"]
                \arrow[d, two heads, "\widehat e"']
                \arrow[dr, phantom, "\lrcorner", very near start]
            &
            P
                \arrow[d, two heads, "e"]
            \\
            A_1\times_{A_0}I
                \arrow[r, "\pi_I"']
            &
            I .
        \end{tikzcd}
        \]

        \item \textbf{Construct the transition by image descent.}
        Consider the composite
        \[
            h_\delta :=
            \delta_Q(1\times m):
            A_1\times_{t_A,A_0,p_I}I\to Q.
        \]
        After precomposition with \(\widehat e\), one has
        \[
        \begin{aligned}
            h_\delta\widehat e(a,x)
            &=
            \delta_Q(a,mex) \\
            &=
            \delta_Q(a,fx) \\
            &=
            f(\delta_P(a,x)) \\
            &=
            me(\delta_P(a,x)).
        \end{aligned}
        \]
        Hence
        \[
            h_\delta\widehat e
        \]
        factors through the underlying monomorphism
        \[
            m:I\hookrightarrow Q.
        \]
        By Lemma~\ref{lem:regular-epi-descent-for-subobjects-membership}, applied in
        \(\mathcal E\), \(h_\delta\) factors uniquely through \(m\). Thus there
        is a unique map
        \[
            \delta_I:
            A_1\times_{t_A,A_0,p_I}I\to I
        \]
        satisfying
        \[
            m\delta_I
            =
            \delta_Q(1\times m).
        \]
        The descent square is
        \[
        \begin{tikzcd}[column sep=large, row sep=large]
            A_1\times_{A_0}P
                \arrow[r, "\widehat e", two heads]
                \arrow[d, "\delta_P"']
            &
            A_1\times_{A_0}I
                \arrow[d, dashed, "\delta_I"]
                \arrow[dr, "{\delta_Q(1\times m)}"]
            \\
            P
                \arrow[r, two heads, "e"']
            &
            I
                \arrow[r, hook, "m"']
            &
            Q .
        \end{tikzcd}
        \]

        \item \textbf{Verify the transition typing.}
        Since \(m\) is a map over \(Z\), the typing equations for \(\delta_I\)
        follow after composing with \(m\):
        \[
        \begin{aligned}
            p_I\delta_I
            &=
            p_Qm\delta_I \\
            &=
            p_Q\delta_Q(1\times m) \\
            &=
            s_A\pi_{A_1},
        \end{aligned}
        \]
        and
        \[
        \begin{aligned}
            q_I\delta_I
            &=
            q_Qm\delta_I \\
            &=
            q_Q\delta_Q(1\times m).
        \end{aligned}
        \]
        The first equation is the left triangle in
        \[
        \begin{tikzcd}[column sep=large, row sep=large]
            A_1\times_{A_0}I
                \arrow[rr, "\delta_I"]
                \arrow[dr, "{s_A\pi_{A_1}}"']
            &&
            I
                \arrow[dl, "p_I"]
                \arrow[r, hook, "m"]
            &
            Q
                \arrow[d, "p_Q"]
            \\
            &
            A_0
                \arrow[rr, equal]
            &&
            A_0 .
        \end{tikzcd}
        \]

        \item \textbf{Define the output by restriction.}
        Define
        \[
            \omega_I :=
            \omega_Q(1\times m):
            A_1\times_{t_A,A_0,p_I}I\to B_1.
        \]
        Its typing equations follow from the typing equations for
        \(\omega_Q\), the equation
        \[
            m\delta_I=\delta_Q(1\times m),
        \]
        and the fact that \(m\) is a map over \(Z\). They are summarized by
        \[
        \begin{tikzcd}[column sep=large, row sep=large]
            A_1\times_{A_0}I
                \arrow[rr, "\omega_I"]
                \arrow[dr, "\delta_I"']
                \arrow[ddr, "\pi_I"', bend right=18]
            &&
            B_1
                \arrow[dl, "s_B"]
                \arrow[ddl, "t_B", bend left=18]
            \\
            &
            I
                \arrow[d, "q_I"]
            &
            \\
            &
            B_0 .
        \end{tikzcd}
        \]

        \item \textbf{Verify strictness of \(m\).}
        The transition equation for \(m\) is the defining equation
        \[
            m\delta_I
            =
            \delta_Q(1\times m).
        \]
        The output equation for \(m\) is the definition
        \[
            \omega_I
            =
            \omega_Q(1\times m).
        \]
        Therefore
        \[
            m:I\hookrightarrow Q
        \]
        is strict.

        \item \textbf{Verify strictness of \(e\).}
        Since
        \[
            me=f,
        \]
        one has
        \[
        \begin{aligned}
            me\delta_P
            &=
            f\delta_P \\
            &=
            \delta_Q(1\times f) \\
            &=
            \delta_Q(1\times m)(1\times e) \\
            &=
            m\delta_I(1\times e).
        \end{aligned}
        \]
        Since \(m\) is monic,
        \[
            e\delta_P=\delta_I(1\times e).
        \]
        Similarly,
        \[
        \begin{aligned}
            \omega_I(1\times e)
            &=
            \omega_Q(1\times m)(1\times e) \\
            &=
            \omega_Q(1\times f) \\
            &=
            \omega_P.
        \end{aligned}
        \]
        Hence
        \[
            e:P\twoheadrightarrow I
        \]
        is strict.

        \item \textbf{Verify the Mealy laws.}
        The unit and multiplication laws for the transition of \(I\) hold after
        composing with \(m\), where they become the corresponding transition
        laws in \(Q\). Since \(m\) is monic, they hold in \(I\).

        The output unit and multiplication laws follow directly from
        \[
            \omega_I=\omega_Q(1\times m)
        \]
        and the output laws in \(Q\), using the transition equation
        \[
            m\delta_I=\delta_Q(1\times m).
        \]

        \item \textbf{Conclude uniqueness.}
        Any Mealy structure on \(I\) making both \(e\) and \(m\) strict must
        have transition and output satisfying the displayed equations. The
        transition is unique by monicity of \(m\), and the output is forced by
        strictness of \(m\). Hence the structure is unique.
    \end{enumerate}
\end{proof}

\subsection{Semantic images}
\label{subsec:semantic-images}

\begin{definition}[Semantic image]
\label{def:semantic-image}
    For a Mealy morphism \(P\), its \textbf{semantic image} is the regular
    image factorization of the unique strict map to the terminal object:
    \[
        P\xrightarrow{e_P}\operatorname{Beh}(P)
        \xrightarrow{m_P}
        \mathbf{Beh}^{\mathrm{can}}_{\mathbb A,\mathbb B}.
    \]
    On carriers, this is the regular image factorization of
    \[
        \beta_P:P\to\mathbf{Beh}_0
    \]
    in \(\mathcal E/Z\).
\end{definition}

The semantic image is displayed by
\[
\begin{tikzcd}[column sep=large, row sep=large]
    P
        \arrow[rr, "\beta_P"]
        \arrow[dr, two heads, "e_P"']
    &&
    \mathbf{Beh}_0
    \\
    &
    \operatorname{Beh}(P)
        \arrow[ur, hook, "m_P"']
    &
\end{tikzcd}
\qquad
\text{in }\mathcal E/Z.
\]
The same factorization, with the Mealy structure created by
Proposition~\ref{prop:regular-image-factorization-mealy}, is
\[
\begin{tikzcd}[column sep=large, row sep=large]
    P
        \arrow[rr, "\beta_P"]
        \arrow[dr, two heads, "e_P"']
    &&
    \mathbf{Beh}^{\mathrm{can}}
    \\
    &
    \operatorname{Beh}(P)
        \arrow[ur, hook, "m_P"']
    &
\end{tikzcd}
\qquad
\text{in }\mathsf{Mly}^{\mathrm{str}}_{\mathbb A,\mathbb B}.
\]

\begin{corollary}[Semantic images are coequational theories]
\label{cor:semantic-images-are-coeq-theories}
    For every Mealy morphism \(P\), the semantic image
    \[
        \operatorname{Beh}(P)\hookrightarrow\mathbf{Beh}^{\mathrm{can}}
    \]
    is a local coequational theory. Moreover,
    \[
        P\twoheadrightarrow\operatorname{Beh}(P)
    \]
    is a regular quotient in
    \[
        \mathsf{Mly}^{\mathrm{str}}_{\mathbb A,\mathbb B}.
    \]
\end{corollary}

\begin{proof}
    \leavevmode
    \begin{enumerate}
        \item \textbf{Apply image creation.}
        The semantic image is the regular image of the strict semantic
        homomorphism
        \[
            \beta_P:P\to\mathbf{Beh}^{\mathrm{can}}.
        \]
        Proposition~\ref{prop:regular-image-factorization-mealy} creates this
        image in the strict Mealy category.

        \item \textbf{Identify the image as a subobject of the terminal object.}
        The monic part
        \[
            \operatorname{Beh}(P)\hookrightarrow\mathbf{Beh}^{\mathrm{can}}
        \]
        is a subobject, hence a local coequational theory.

        \item \textbf{Identify the quotient map.}
        The regular-epimorphic part
        \[
            P\twoheadrightarrow\operatorname{Beh}(P)
        \]
        is strict by the same proposition.
    \end{enumerate}
\end{proof}

\subsection{Meets and coproducts}
\label{subsec:meets-and-coproducts-local-theories}

\begin{lemma}[Meets of local theories]
\label{lem:meets-local-theories}
    If
    \[
        S,T\hookrightarrow\mathbf{Beh}^{\mathrm{can}}
    \]
    are local coequational theories, then their meet
    \[
        S\wedge T\hookrightarrow\mathbf{Beh}^{\mathrm{can}}
    \]
    is a local coequational theory and is the pullback subobject.
\end{lemma}

\begin{proof}
    \leavevmode
    \begin{enumerate}
        \item \textbf{Take the pullback of subobjects.}
        The meet is the pullback in the carrier slice:
        \[
        \begin{tikzcd}[column sep=large, row sep=large]
            S\wedge T
                \arrow[r]
                \arrow[d]
                \arrow[dr, phantom, "\lrcorner", very near start]
            &
            T
                \arrow[d, hook]
            \\
            S
                \arrow[r, hook]
            &
            \mathbf{Beh}_0 .
        \end{tikzcd}
        \]

        \item \textbf{Restrict the Mealy structure.}
        Since both \(S\) and \(T\) are subobjects of
        \(\mathbf{Beh}^{\mathrm{can}}\), the transition and output of the
        canonical object restrict to their intersection. The transition
        restriction is the pullback comparison
        \[
        \begin{tikzcd}[column sep=huge, row sep=large]
            A_1\times_{A_0}(S\wedge T)
                \arrow[r]
                \arrow[d, dashed, "\delta_{S\wedge T}"']
            &
            A_1\times_{A_0}T
                \arrow[d, "\delta_T"]
            \\
            S\wedge T
                \arrow[r]
                \arrow[d]
            &
            T
                \arrow[d, hook]
            \\
            S
                \arrow[r, hook]
            &
            \mathbf{Beh}_0 .
        \end{tikzcd}
        \]
        The laws hold by monicity.

        \item \textbf{Conclude the universal property.}
        The resulting object is the pullback in
        \(\mathsf{Mly}^{\mathrm{str}}_{\mathbb A,\mathbb B}\), and its carrier
        is the meet in \(\mathcal E/Z\).
    \end{enumerate}
\end{proof}

\begin{proposition}[Small coproducts are created by the carrier slice]
\label{prop:small-coproducts-mealy}
    A small family
    \[
        (P_i)_{i\in I}
    \]
    of Mealy morphisms over \((\mathbb A,\mathbb B)\) has a coproduct in
    \[
        \mathsf{Mly}^{\mathrm{str}}_{\mathbb A,\mathbb B}.
    \]
    Its carrier is the coproduct
    \[
        \coprod_{i\in I}P_i
    \]
    in the slice
    \[
        \mathcal E/Z,
    \]
    and its transition and output maps are induced componentwise.
\end{proposition}

\begin{proof}
    \leavevmode
    \begin{enumerate}
        \item \textbf{Use preservation of coproducts by pullback.}
        In a Grothendieck topos, pullback preserves small coproducts in every
        slice. Hence
        \[
            A_1\times_{t_A,A_0,p_{\coprod_iP_i}}
            \coprod_iP_i
            \cong
            \coprod_i
            \left(
                A_1\times_{t_A,A_0,p_{P_i}}P_i
            \right).
        \]
        The coproduct of transition domains is displayed by
        \[
        \begin{tikzcd}[column sep=large, row sep=large]
            A_1\times_{A_0}P_i
                \arrow[r]
                \arrow[d]
            &
            A_1\times_{A_0}\coprod_iP_i
                \arrow[d]
            \\
            P_i
                \arrow[r]
            &
            \coprod_iP_i .
        \end{tikzcd}
        \]

        \item \textbf{Define transition and output componentwise.}
        The transition and output maps are induced by the component maps
        \[
            \delta_{P_i},
            \qquad
            \omega_{P_i}.
        \]
        Equivalently, for each \(i\) the following diagrams commute:
        \[
        \begin{tikzcd}[column sep=huge, row sep=large]
            A_1\times_{A_0}P_i
                \arrow[r]
                \arrow[d, "\delta_{P_i}"']
            &
            A_1\times_{A_0}\coprod_iP_i
                \arrow[d, "\delta_{\coprod_iP_i}"]
            \\
            P_i
                \arrow[r]
            &
            \coprod_iP_i
        \end{tikzcd}
        \qquad
        \begin{tikzcd}[column sep=huge, row sep=large]
            A_1\times_{A_0}P_i
                \arrow[r]
                \arrow[dr, "\omega_{P_i}"']
            &
            A_1\times_{A_0}\coprod_iP_i
                \arrow[d, "\omega_{\coprod_iP_i}"]
            \\
            &
            B_1 .
        \end{tikzcd}
        \]

        \item \textbf{Verify the laws.}
        The Mealy laws hold after restriction to each coproduct summand. Since
        the summand inclusions are jointly epic, the laws hold on the
        coproduct.

        \item \textbf{Check the universal property.}
        A strict semantic homomorphism
        \[
            \coprod_iP_i\to Q
        \]
        is equivalently a compatible family of strict semantic homomorphisms
        \[
            P_i\to Q.
        \]
        This is the universal property of the coproduct in \(\mathcal E/Z\)
        together with the componentwise Mealy structure:
        \[
        \begin{tikzcd}[column sep=large, row sep=large]
            P_i
                \arrow[r]
                \arrow[dr, "f_i"']
            &
            \coprod_iP_i
                \arrow[d, dashed, "f"]
            \\
            &
            Q .
        \end{tikzcd}
        \]
    \end{enumerate}
\end{proof}

\subsection{Running specializations}
\label{subsec:semantic-images-running-specializations}

In the one-object Boolean language case, the semantic image of a recognizer is
the set of all residual languages generated by its states:
\[
    \operatorname{Beh}(P)
    =
    \{\beta_P(x)\mid x\in P\}
    \subseteq
    2^{\Sigma^*}.
\]
The quotient
\[
    P\twoheadrightarrow\operatorname{Beh}(P)
\]
identifies states with the same residual language. This is the semantic
quotient underlying the classical Nerode quotient.

In the stateless case, for an internal functor
\[
    F:\mathbb A\to\mathbb B,
\]
the semantic image is the regular image of
\[
    A_0\to\mathbf{Beh}_0,
    \qquad
    v\longmapsto H^F_v.
\]
Thus it records the residual slice functors generated by \(F\), modulo equality
in the behaviour object.

\section{The local coequational Birkhoff theorem}
\label{sec:local-coequational-birkhoff}

\subsection{Closure properties of model classes}
\label{subsec:closure-properties-model-classes}

\begin{theorem}[Local coequational closure]
\label{thm:local-coequational-closure}
    Let
    \[
        V\hookrightarrow\mathbf{Beh}^{\mathrm{can}}
    \]
    be a local coequational theory. Then:
    \begin{enumerate}
        \item \(\mathsf{Mod}(V)\) is closed under strict homomorphic preimages.
        \item \(\mathsf{Mod}(V)\) is closed under regular-epimorphic quotients.
        \item \(\mathsf{Mod}(V)\) is closed under small coproducts.
        \item \(V\), regarded as the restricted canonical object, is itself a
        model of \(V\), and
        \[
            \operatorname{Beh}(V)=V.
        \]
    \end{enumerate}
    Consequently the assignment
    \[
        V\longmapsto\mathsf{Mod}(V)
    \]
    is monotone and order-reflecting.
\end{theorem}

\begin{proof}
    Let
    \[
        j_V:V\hookrightarrow\mathbf{Beh}^{\mathrm{can}}
    \]
    denote the inclusion.

    \begin{enumerate}
        \item \textbf{Closure under strict homomorphic preimages.}
        Let
        \[
            f:P\to Q
        \]
        be strict and suppose \(Q\models V\). Choose
        \[
            \bar\beta_Q:Q\to V
        \]
        with
        \[
            j_V\bar\beta_Q=\beta_Q.
        \]
        Then, by Lemma~\ref{lem:coeq-strict-maps-preserve-beta},
        \[
            \beta_P=\beta_Qf=j_V\bar\beta_Qf,
        \]
        so \(P\models V\). The factorization is
        \[
        \begin{tikzcd}[column sep=large, row sep=large]
            P
                \arrow[r, "f"]
                \arrow[rr, bend left=16, "\beta_P"]
                \arrow[dr, dashed, "\bar\beta_Qf"']
            &
            Q
                \arrow[r, "\beta_Q"]
                \arrow[d, "\bar\beta_Q"]
            &
            \mathbf{Beh}^{\mathrm{can}}
            \\
            &
            V
                \arrow[ur, hook, "j_V"']
            &
        \end{tikzcd}
        \]

        \item \textbf{Closure under regular-epimorphic quotients.}
        Let
        \[
            e:P\twoheadrightarrow Q
        \]
        be a regular-epimorphic strict semantic homomorphism and suppose
        \(P\models V\). Since
        \[
            \beta_P=\beta_Qe,
        \]
        and \(\beta_P\) factors through
        \[
            j_V:V\hookrightarrow\mathbf{Beh}^{\mathrm{can}},
        \]
        Lemma~\ref{lem:regular-epi-descent-for-subobjects-membership} implies that
        \(\beta_Q\) factors through \(j_V\). Hence \(Q\models V\).
        Diagrammatically,
        \[
        \begin{tikzcd}[column sep=large, row sep=large]
            P
                \arrow[r, two heads, "e"]
                \arrow[dr, "\bar\beta_P"']
                \arrow[ddr, "\beta_P"', bend right=18]
            &
            Q
                \arrow[d, dashed, "\bar\beta_Q"]
                \arrow[dd, "\beta_Q", bend left=18]
            \\
            &
            V
                \arrow[d, hook, "j_V"]
            \\
            &
            \mathbf{Beh}^{\mathrm{can}} .
        \end{tikzcd}
        \]

        \item \textbf{Closure under small coproducts.}
        Let
        \[
            (P_i)_{i\in I}
        \]
        be a small family of models. For each \(i\), the factorization
        \[
            \bar\beta_i:P_i\to V
        \]
        is strict by Lemma~\ref{lem:model-factorizations-strict}. By the
        coproduct universal property in
        \(\mathsf{Mly}^{\mathrm{str}}_{\mathbb A,\mathbb B}\), these strict
        maps induce a strict map
        \[
            \coprod_iP_i\to V.
        \]
        Therefore the coproduct models \(V\). The construction is
        \[
        \begin{tikzcd}[column sep=large, row sep=large]
            P_i
                \arrow[r]
                \arrow[dr, "\bar\beta_i"']
            &
            \coprod_iP_i
                \arrow[d, dashed, "\coprod_i\bar\beta_i"]
                \arrow[dd, bend left=18, "\beta_{\coprod_iP_i}"]
            \\
            &
            V
                \arrow[d, hook, "j_V"]
            \\
            &
            \mathbf{Beh}^{\mathrm{can}} .
        \end{tikzcd}
        \]

        \item \textbf{The theory models itself.}
        The identity map
        \[
            V\to V
        \]
        exhibits \(V\) as a model of itself. By
        Lemma~\ref{lem:semantic-map-of-a-restricted-canonical-subobject}, the
        semantic map of a restricted canonical subobject is its inclusion
        \[
            j_V:V\hookrightarrow\mathbf{Beh}^{\mathrm{can}}.
        \]
        Hence
        \[
            \operatorname{Beh}(V)=V.
        \]

        \item \textbf{Prove monotonicity and order-reflection.}
        If \(V\leq W\), then every factorization through \(V\) is a
        factorization through \(W\), so
        \[
            \mathsf{Mod}(V)\subseteq\mathsf{Mod}(W).
        \]
        Conversely, if
        \[
            \mathsf{Mod}(V)\subseteq\mathsf{Mod}(W),
        \]
        then \(V\in\mathsf{Mod}(V)\), hence \(V\in\mathsf{Mod}(W)\). Therefore
        the inclusion
        \[
            V\hookrightarrow\mathbf{Beh}^{\mathrm{can}}
        \]
        factors through
        \[
            W\hookrightarrow\mathbf{Beh}^{\mathrm{can}},
        \]
        so \(V\leq W\). This is the triangle
        \[
        \begin{tikzcd}[column sep=large, row sep=large]
            V
                \arrow[rr, hook, "j_V"]
                \arrow[dr, dashed]
            &&
            \mathbf{Beh}^{\mathrm{can}}
            \\
            &
            W
                \arrow[ur, hook, "j_W"']
            &
        \end{tikzcd}
        \]
        obtained from the fact that \(V\) is a model of \(W\).
    \end{enumerate}
\end{proof}

Closure under small coproducts includes the empty coproduct. Hence the empty
replete full subcategory is not coequationally definable.

\subsection{Joins of local theories}
\label{subsec:joins-local-theories}

Let
\[
    \mathcal C\subseteq\mathsf{Mly}^{\mathrm{str}}_{\mathbb A,\mathbb B}
\]
be a replete full subcategory. Define
\[
    \mathcal S_{\mathcal C}
    :=
    \{
        S\hookrightarrow\mathbf{Beh}^{\mathrm{can}}
        \mid
        S=\operatorname{Beh}(P)
        \text{ for some }P\in\mathcal C
    \}.
\]
Although \(\mathcal C\) may be large, the collection
\[
    \mathcal S_{\mathcal C}
\]
is a subclass of the set
\[
    \operatorname{Sub}_{\mathcal E/Z}(\mathbf{Beh}_0).
\]
We therefore choose a small set of representatives for the subobjects in
\(\mathcal S_{\mathcal C}\) and form their join. Define
\[
    V_{\mathcal C}
    :=
    \bigvee_{S\in\mathcal S_{\mathcal C}}S
    \hookrightarrow\mathbf{Beh}^{\mathrm{can}},
\]
where the displayed join is taken over the chosen small representative set.

\begin{lemma}[Joins of local coequational theories]
\label{lem:joins-derivative-closed}
    Let
    \[
        (S_i\hookrightarrow\mathbf{Beh}^{\mathrm{can}})_{i\in I}
    \]
    be a small family of local coequational theories. Then their join
    \[
        S:=\bigvee_{i\in I}S_i
    \]
    in
    \[
        \operatorname{Sub}_{\mathcal E/Z}(\mathbf{Beh}_0)
    \]
    is again a local coequational theory.
\end{lemma}

\begin{proof}
    \leavevmode
    \begin{enumerate}
        \item \textbf{Represent the join.}
        In the Grothendieck topos \(\mathcal E/Z\), the join is represented by
        the regular image of the coproduct map
        \[
            \coprod_iS_i\to\mathbf{Beh}_0.
        \]
        Write this image factorization as
        \[
        \begin{tikzcd}[column sep=large, row sep=large]
            \coprod_iS_i
                \arrow[rr]
                \arrow[dr, two heads]
            &&
            \mathbf{Beh}_0
            \\
            &
            S
                \arrow[ur, hook]
            &
        \end{tikzcd}
        \]
        in \(\mathcal E/Z\).

        \item \textbf{Regard the coproduct as a Mealy object.}
        By Proposition~\ref{prop:small-coproducts-mealy}, the coproduct
        \[
            \coprod_iS_i
        \]
        carries the componentwise Mealy structure. Since each
        \[
            S_i\hookrightarrow\mathbf{Beh}^{\mathrm{can}}
        \]
        is strict, the induced coproduct map to
        \(\mathbf{Beh}^{\mathrm{can}}\) is strict. Strictness is tested after
        precomposition with the jointly epic coproduct injections:
        \[
        \begin{tikzcd}[column sep=large, row sep=large]
            S_i
                \arrow[r]
                \arrow[dr, hook]
            &
            \coprod_iS_i
                \arrow[d, dashed]
            \\
            &
            \mathbf{Beh}^{\mathrm{can}} .
        \end{tikzcd}
        \]

        \item \textbf{Take the regular image in the strict Mealy category.}
        By Proposition~\ref{prop:regular-image-factorization-mealy}, the
        regular image of this strict map is created in \(\mathcal E/Z\) and
        carries a Mealy structure. Hence
        \[
            S\hookrightarrow\mathbf{Beh}^{\mathrm{can}}
        \]
        is a subobject in the strict Mealy category. The created image is
        \[
        \begin{tikzcd}[column sep=large, row sep=large]
            \coprod_iS_i
                \arrow[rr]
                \arrow[dr, two heads]
            &&
            \mathbf{Beh}^{\mathrm{can}}
            \\
            &
            S
                \arrow[ur, hook]
            &
        \end{tikzcd}
        \]
        in \(\mathsf{Mly}^{\mathrm{str}}_{\mathbb A,\mathbb B}\).

        \item \textbf{Conclude local coequationality.}
        By Proposition~\ref{prop:subobjects-canonical-behaviour-machine},
        being a subobject of the canonical behaviour object is precisely being
        a local coequational theory.
    \end{enumerate}
\end{proof}

\subsection{The theorem}
\label{subsec:local-coequational-birkhoff-theorem}

\begin{theorem}[Local coequational Birkhoff theorem]
\label{thm:local-coequational-birkhoff}
    Let
    \[
        \mathcal C\subseteq\mathsf{Mly}^{\mathrm{str}}_{\mathbb A,\mathbb B}
    \]
    be a replete full subcategory. Then \(\mathcal C\) is of the form
    \[
        \mathcal C=\mathsf{Mod}(V)
    \]
    for a local coequational theory
    \[
        V\hookrightarrow\mathbf{Beh}^{\mathrm{can}}_{\mathbb A,\mathbb B}
    \]
    if and only if \(\mathcal C\) is closed under:
    \begin{enumerate}
        \item strict homomorphic preimages;
        \item regular-epimorphic quotients;
        \item small coproducts.
    \end{enumerate}
    When these conditions hold, the defining theory is
    \[
        V_{\mathcal C}
        =
        \bigvee_{P\in\mathcal C}\operatorname{Beh}(P)
        \hookrightarrow\mathbf{Beh}^{\mathrm{can}},
    \]
    where the join is taken over a chosen small set of representatives for the
    resulting semantic-image subobjects.
\end{theorem}

\begin{proof}
    \leavevmode
    \begin{enumerate}
        \item \textbf{Prove necessity.}
        If
        \[
            \mathcal C=\mathsf{Mod}(V),
        \]
        then the desired closure properties are
        Theorem~\ref{thm:local-coequational-closure}.

        \item \textbf{Construct the candidate theory.}
        Conversely, suppose \(\mathcal C\) is closed under strict homomorphic
        preimages, regular-epimorphic quotients, and small coproducts. Define
        \[
            V_{\mathcal C}
            :=
            \bigvee_{P\in\mathcal C}\operatorname{Beh}(P)
            \hookrightarrow\mathbf{Beh}^{\mathrm{can}},
        \]
        taking the join over representatives of the distinct semantic-image
        subobjects. Each \(\operatorname{Beh}(P)\) is a local coequational
        theory by
        Corollary~\ref{cor:semantic-images-are-coeq-theories}, so
        \(V_{\mathcal C}\) is a local coequational theory by
        Lemma~\ref{lem:joins-derivative-closed}.

        \item \textbf{Show \(\mathcal C\subseteq\mathsf{Mod}(V_{\mathcal C})\).}
        If \(P\in\mathcal C\), then
        \[
            \operatorname{Beh}(P)\leq V_{\mathcal C}.
        \]
        Since \(\beta_P\) factors through \(\operatorname{Beh}(P)\), it factors
        through \(V_{\mathcal C}\). Hence \(P\models V_{\mathcal C}\):
        \[
        \begin{tikzcd}[column sep=large, row sep=large]
            P
                \arrow[r, two heads]
                \arrow[rrr, bend left=18, "\beta_P"]
            &
            \operatorname{Beh}(P)
                \arrow[r, hook]
            &
            V_{\mathcal C}
                \arrow[r, hook]
            &
            \mathbf{Beh}^{\mathrm{can}} .
        \end{tikzcd}
        \]

        \item \textbf{Let a model be given.}
        Let
        \[
            X\in\mathsf{Mod}(V_{\mathcal C}).
        \]
        Then
        \[
            \operatorname{Beh}(X)\leq V_{\mathcal C}.
        \]

        \item \textbf{Decompose the semantic image by intersections.}
        Since \(\mathcal E/Z\) is again a Grothendieck topos, the lattice
        \[
            \operatorname{Sub}_{\mathcal E/Z}(\mathbf{Beh}_0)
        \]
        is a frame: finite meets distribute over small joins. Therefore
        \[
        \begin{aligned}
            \operatorname{Beh}(X)
            &=
            \operatorname{Beh}(X)\wedge V_{\mathcal C} \\
            &=
            \operatorname{Beh}(X)\wedge
            \bigvee_{S\in\mathcal S_{\mathcal C}}S \\
            &=
            \bigvee_{S\in\mathcal S_{\mathcal C}}
            \bigl(\operatorname{Beh}(X)\wedge S\bigr).
        \end{aligned}
        \]
        The intersections are the pullbacks
        \[
        \begin{tikzcd}[column sep=large, row sep=large]
            \operatorname{Beh}(X)\wedge S
                \arrow[r]
                \arrow[d]
                \arrow[dr, phantom, "\lrcorner", very near start]
            &
            S
                \arrow[d, hook]
            \\
            \operatorname{Beh}(X)
                \arrow[r, hook]
            &
            \mathbf{Beh}^{\mathrm{can}} .
        \end{tikzcd}
        \]

        \item \textbf{Show each intersection belongs to \(\mathcal C\).}
        Choose \(P_S\in\mathcal C\) with
        \[
            S=\operatorname{Beh}(P_S).
        \]
        Since
        \[
            P_S\twoheadrightarrow S
        \]
        is a regular quotient and \(\mathcal C\) is closed under regular
        quotients, we have
        \[
            S\in\mathcal C,
        \]
        using repleteness if necessary to replace \(S\) by the chosen
        representative. By Lemma~\ref{lem:meets-local-theories},
        \[
            \operatorname{Beh}(X)\wedge S
        \]
        is a local coequational theory and the projection
        \[
            \operatorname{Beh}(X)\wedge S\hookrightarrow S
        \]
        is a monic strict semantic homomorphism. Closure under strict
        homomorphic preimages gives
        \[
            \operatorname{Beh}(X)\wedge S\in\mathcal C.
        \]

        \item \textbf{Use coproducts and regular images.}
        Form
        \[
            C_X:=
            \coprod_{S\in\mathcal S_{\mathcal C}}
            \bigl(\operatorname{Beh}(X)\wedge S\bigr),
        \]
        using the chosen representative set for
        \(\mathcal S_{\mathcal C}\). Each summand lies in \(\mathcal C\), and
        \(\mathcal C\) is closed under small coproducts, so
        \[
            C_X\in\mathcal C.
        \]
        The canonical strict map
        \[
            C_X\to\mathbf{Beh}^{\mathrm{can}}
        \]
        has regular image
        \[
            \bigvee_{S\in\mathcal S_{\mathcal C}}
            \bigl(\operatorname{Beh}(X)\wedge S\bigr)
            =
            \operatorname{Beh}(X).
        \]
        The image diagram is
        \[
        \begin{tikzcd}[column sep=large, row sep=large]
            \coprod_S\bigl(\operatorname{Beh}(X)\wedge S\bigr)
                \arrow[rr]
                \arrow[dr, two heads]
            &&
            \mathbf{Beh}^{\mathrm{can}}
            \\
            &
            \operatorname{Beh}(X)
                \arrow[ur, hook]
            &
        \end{tikzcd}
        \]
        in \(\mathsf{Mly}^{\mathrm{str}}_{\mathbb A,\mathbb B}\).
        By Proposition~\ref{prop:regular-image-factorization-mealy}, this image
        is a regular quotient of \(C_X\). Since \(\mathcal C\) is closed under
        regular-epimorphic quotients,
        \[
            \operatorname{Beh}(X)\in\mathcal C.
        \]

        \item \textbf{Pull back along the semantic quotient.}
        The semantic image factorization gives a regular quotient
        \[
            X\twoheadrightarrow\operatorname{Beh}(X).
        \]
        Since \(\operatorname{Beh}(X)\in\mathcal C\) and \(\mathcal C\) is
        closed under strict homomorphic preimages, it follows that
        \[
            X\in\mathcal C.
        \]
        This is the final use of
        \[
        \begin{tikzcd}[column sep=large, row sep=large]
            X
                \arrow[r, two heads]
                \arrow[rr, bend left=16, "\beta_X"]
            &
            \operatorname{Beh}(X)
                \arrow[r, hook]
            &
            \mathbf{Beh}^{\mathrm{can}} .
        \end{tikzcd}
        \]

        \item \textbf{Conclude equality.}
        Thus
        \[
            \mathsf{Mod}(V_{\mathcal C})\subseteq\mathcal C.
        \]
        Together with the opposite inclusion, this gives
        \[
            \mathcal C=\mathsf{Mod}(V_{\mathcal C}).
        \]
    \end{enumerate}
\end{proof}

\subsection{Running specializations}
\label{subsec:local-birkhoff-running-specializations}

In the one-object Boolean language case, the local theorem says that
derivative-closed collections of languages correspond to classes of recognizers
closed under semantic preimages, quotients, and coproducts. The representing
theory of a class \(\mathcal C\) is
\[
    V_{\mathcal C}
    =
    \bigcup_{P\in\mathcal C}\operatorname{Beh}(P)
    \subseteq
    2^{\Sigma^*},
\]
closed under derivatives by construction. In the topos-theoretic statement,
this union is the join of subobjects.

In the stateless case, the theorem says that a class of internal functors
viewed through their stateless Mealy morphisms is determined by the residual
slice functors it generates, provided the class is closed under the strict
preimage, quotient, and coproduct operations available in
\[
    \mathsf{Mly}^{\mathrm{str}}_{\mathbb A,\mathbb B}.
\]

\section{Structured output and algebraic closure}
\label{sec:structured-output-algebraic-closure}

\subsection{Structured output categories}
\label{subsec:structured-output-categories}

This section explains how additional closure operations on behaviours arise
from algebraic structure on the output category. The categorical treatment of
varieties of languages uses such algebraic closure conditions as part of the
bridge between syntactic structure and semantic classes
\cite{Eilenberg1974,AdamekMiliusMyersUrbat2019TOCL,Salamanca2016Monad}.

\begin{definition}[Structured output category]
\label{def:structured-output-category}
    Let \(\mathbb T\) be a single-sorted finite-product theory. A
    \textbf{\(\mathbb T\)-structured output category} is a \(\mathbb T\)-algebra
    object in internal categories, equivalently a finite-product-preserving
    functor
    \[
        \mathbb T\longrightarrow \mathbf{Cat}_{\mathrm{int}}(\mathcal E).
    \]
    We write
    \[
        \mathbb B
    \]
    for the internal category assigned to the generic object of \(\mathbb T\).
    Thus each operation of arity \(n\) in \(\mathbb T\) determines an internal
    functor
    \[
        \theta:\mathbb B^n\to\mathbb B,
    \]
    and the equations of \(\mathbb T\) hold internally.
\end{definition}

\subsection{Pointwise operations on behaviours}
\label{subsec:pointwise-operations-behaviours}

For \(n>0\), put
\[
    \mathbf{Beh}_0^{(n)}
    :=
    \underbrace{
    \mathbf{Beh}_0\times_{A_0}\cdots\times_{A_0}\mathbf{Beh}_0
    }_{n}.
\]
A generalized element is a tuple
\[
    (H_1,\dots,H_n)
\]
of residual behaviours
\[
    H_i:\mathbb A/v\to\mathbb B
\]
over the same input object \(v\). For \(n=0\), set
\[
    \mathbf{Beh}_0^{(0)}:=A_0.
\]

For an \(n\)-ary operation
\[
    \theta:\mathbb B^n\to\mathbb B,
\]
we also use the notation
\[
    \mathbf{Beh}^{(n)}_{\theta,0}
\]
for the same underlying object \(\mathbf{Beh}_0^{(n)}\), equipped with its
\(\theta\)-induced map to
\[
    Z=A_0\times B_0.
\]
For \(n>0\), this \(Z\)-coordinate is
\[
    (H_1,\dots,H_n)
    \longmapsto
    \Bigl(
        p_{\mathbf{Beh}}(H_1),
        \theta_0(q_{\mathbf{Beh}}H_1,\dots,q_{\mathbf{Beh}}H_n)
    \Bigr).
\]
This coordinate is displayed by
\[
\begin{tikzcd}[column sep=large, row sep=large]
    \mathbf{Beh}^{(n)}_{\theta,0}
        \arrow[r, "{(q_{\mathbf{Beh}}\pi_i)_i}"]
        \arrow[d, "{p_{\mathbf{Beh}}\pi_1}"']
        \arrow[dr, dashed, "{(p,\theta_0(q_i)_i)}"]
    &
    B_0^n
        \arrow[d, "\theta_0"]
    \\
    A_0
        \arrow[r]
    &
    Z=A_0\times B_0 .
\end{tikzcd}
\]
For \(n=0\), if \(\theta:1\to\mathbb B\) selects the object
\[
    b_\theta:1\to B_0,
\]
then
\[
    \mathbf{Beh}^{(0)}_{\theta,0}=A_0
\]
has \(Z\)-coordinate
\[
    v\longmapsto(v,b_\theta).
\]

\begin{definition}[Pointwise structured behaviour]
\label{def:pointwise-structured-behaviour}
    Let
    \[
        \theta:\mathbb B^n\to\mathbb B
    \]
    be an operation of \(\mathbb T\). The induced map
    \[
        \theta_{\mathbf{Beh}}:
        \mathbf{Beh}^{(n)}_{\theta,0}\to\mathbf{Beh}_0
    \]
    in \(\mathcal E/Z\) is defined by
    \[
        \theta_{\mathbf{Beh}}(H_1,\dots,H_n)
        =
        \theta\circ\langle H_1,\dots,H_n\rangle
    \]
    for \(n>0\). For \(n=0\), the nullary operation
    \[
        \theta:1\to\mathbb B
    \]
    induces the constant residual behaviour
    \[
        \mathbb A/v\to 1\xrightarrow{\theta}\mathbb B
    \]
    over each \(v\in A_0\).
\end{definition}

The construction is summarized by
\[
\begin{tikzcd}[column sep=huge, row sep=large]
    \mathbb A/v
        \arrow[r, "{\langle H_1,\dots,H_n\rangle}"]
        \arrow[rr, "{\theta_{\mathbf{Beh}}(H_i)_i}"', bend right=14]
    &
    \mathbb B^n
        \arrow[r, "\theta"]
    &
    \mathbb B .
\end{tikzcd}
\]
On anchor objects and residual arrows this is the pair of diagrams
\[
\begin{tikzcd}[column sep=large, row sep=large]
    1_v
        \arrow[r, maps to]
        \arrow[d, maps to]
    &
    (H_1(1_v),\dots,H_n(1_v))
        \arrow[d, "\theta_0"]
    \\
    \theta_{\mathbf{Beh}}(H_i)_i(1_v)
        \arrow[r, equal]
    &
    \theta_0(H_1(1_v),\dots,H_n(1_v))
\end{tikzcd}
\]
and
\[
\begin{tikzcd}[column sep=large, row sep=large]
    r\circ a
        \arrow[r]
        \arrow[d, maps to]
    &
    r
        \arrow[d, maps to]
    \\
    (H_i(r\circ a))_i
        \arrow[r, "{(H_i(r\circ a\to r))_i}"']
        \arrow[d, "\theta_0"']
    &
    (H_i(r))_i
        \arrow[d, "\theta_0"]
    \\
    \theta(H_i(r\circ a))_i
        \arrow[r, "{\theta_1(H_i(r\circ a\to r))_i}"']
    &
    \theta(H_i(r))_i .
\end{tikzcd}
\]

For a nullary operation \(\theta:1\to\mathbb B\), write
\[
    \theta()
\]
for the Mealy morphism with carrier \(A_0\), input projection \(1_{A_0}\),
constant output \(b_\theta\), transition
\[
    \delta_{\theta()}(a,v)=s_A(a),
\]
and identity output arrow at \(b_\theta\).

\begin{lemma}[Structured operations on strict Mealy morphisms]
\label{lem:structured-operations-semantic}
    For every operation
    \[
        \theta:\mathbb B^n\to\mathbb B,
    \]
    there is an induced \(n\)-ary functor
    \[
        \theta:
        \bigl(\mathsf{Mly}^{\mathrm{str}}_{\mathbb A,\mathbb B}\bigr)^n
        \to
        \mathsf{Mly}^{\mathrm{str}}_{\mathbb A,\mathbb B}
    \]
    for \(n>0\), with the evident nullary object for \(n=0\). The residual
    semantics is compatible with this operation:
    \[
        \beta_{\theta(P_i)_i}
        =
        \theta_{\mathbf{Beh}}(\beta_{P_i})_i.
    \]
\end{lemma}

\begin{proof}
    \leavevmode
    \begin{enumerate}
        \item \textbf{Construct the carrier.}
        For \(n>0\), the carrier is
        \[
            P_\theta:=P_1\times_{A_0}\cdots\times_{A_0}P_n.
        \]
        Its input projection is the common projection to \(A_0\), and its
        output projection is
        \[
            q_\theta(x_1,\dots,x_n)
            =
            \theta_0(q_1x_1,\dots,q_nx_n).
        \]
        This is the pullback-and-output diagram
        \[
        \begin{tikzcd}[column sep=large, row sep=large]
            P_\theta
                \arrow[r, "\pi_i"]
                \arrow[d, "p_\theta"']
                \arrow[dr, phantom, "\lrcorner", very near start]
                \arrow[rr, bend left=18, "{(q_i\pi_i)_i}"]
            &
            P_i
                \arrow[d, "p_i"]
                \arrow[r, "q_i"]
            &
            B_0
            \\
            A_0
                \arrow[r, equal]
            &
            A_0
                \arrow[ur, phantom, "\cdots" description]
            &
        \end{tikzcd}
        \]
        together with
        \[
            q_\theta=\theta_0(q_i\pi_i)_i.
        \]
        For \(n=0\), the carrier is \(A_0\), with input projection
        \(1_{A_0}\) and output projection
        \[
            A_0\to 1\xrightarrow{b_\theta}B_0.
        \]

        \item \textbf{Identify the transition domain.}
        For \(n>0\), there is a canonical pullback identification
        \[
        \begin{aligned}
            A_1\times_{A_0}
            \left(P_1\times_{A_0}\cdots\times_{A_0}P_n\right)
            \cong
            \left(A_1\times_{A_0}P_1\right)
            \times_{A_1}
            \cdots
            \times_{A_1}
            \left(A_1\times_{A_0}P_n\right).
        \end{aligned}
        \]
        Under this identification, a transition datum consists of one input
        arrow \(a\) and compatible states \((x_i)_i\). The compatibility is
        expressed by
        \[
        \begin{tikzcd}[column sep=large, row sep=large]
            A_1\times_{A_0}P_\theta
                \arrow[r]
                \arrow[d]
                \arrow[dr, phantom, "\lrcorner", very near start]
            &
            P_\theta
                \arrow[d]
            \\
            A_1
                \arrow[r, "t_A"']
            &
            A_0 .
        \end{tikzcd}
        \]

        \item \textbf{Define transition and output.}
        For \(n>0\), define
        \[
            \delta_\theta(a,(x_i)_i)
            =
            (\delta_i(a,x_i))_i.
        \]
        The output arrows
        \[
            \omega_i(a,x_i):
            q_i(\delta_i(a,x_i))\to q_i(x_i)
        \]
        form an arrow in \(\mathbb B^n\). Applying the internal functor
        \[
            \theta:\mathbb B^n\to\mathbb B
        \]
        gives a well-typed arrow
        \[
            \theta_1(\omega_i(a,x_i))_i:
            \theta_0(q_i\delta_i(a,x_i))_i
            \to
            \theta_0(q_ix_i)_i.
        \]
        Define
        \[
            \omega_\theta(a,(x_i)_i)
            =
            \theta_1(\omega_i(a,x_i))_i.
        \]
        The construction is summarized by
        \[
        \begin{tikzcd}[column sep=large, row sep=large]
            \theta_0(q_i\delta_i(a,x_i))_i
                \arrow[r, "{\theta_1(\omega_i(a,x_i))_i}"]
                \arrow[d, equal]
            &
            \theta_0(q_ix_i)_i
                \arrow[d, equal]
            \\
            q_\theta\delta_\theta(a,(x_i)_i)
                \arrow[r, "{\omega_\theta(a,(x_i)_i)}"']
            &
            q_\theta(x_i)_i .
        \end{tikzcd}
        \]

        For \(n=0\), the transition is
        \[
            \delta_\theta(a,v)=s_A(a),
        \]
        and the output map is the composite
        \[
            A_1
            \xrightarrow{!}
            1
            \xrightarrow{b_\theta}
            B_0
            \xrightarrow{e_B}
            B_1.
        \]
        In generalized-element notation,
        \[
            \omega_\theta(a,v)=e_B(b_\theta).
        \]

        \item \textbf{Define the operation on strict maps.}
        If
        \[
            f_i:P_i\to Q_i
        \]
        are strict semantic homomorphisms, define
        \[
            \theta(f_i)_i:
            P_1\times_{A_0}\cdots\times_{A_0}P_n
            \to
            Q_1\times_{A_0}\cdots\times_{A_0}Q_n
        \]
        componentwise. Strictness follows from the componentwise transition
        equations and from functoriality of
        \[
            \theta:\mathbb B^n\to\mathbb B.
        \]
        The output square is
        \[
        \begin{tikzcd}[column sep=large, row sep=large]
            A_1\times_{A_0}P_\theta
                \arrow[r, "{1\times\theta(f_i)}"]
                \arrow[d, "\omega_\theta^P"']
            &
            A_1\times_{A_0}Q_\theta
                \arrow[d, "\omega_\theta^Q"]
            \\
            B_1
                \arrow[r, equal]
            &
            B_1 .
        \end{tikzcd}
        \]
        Identities and composition are inherited from the product in the
        carrier slice.

        \item \textbf{Verify the Mealy laws structurally.}
        The transition laws are inherited componentwise. The output laws follow
        from functoriality of
        \[
            \theta:\mathbb B^n\to\mathbb B.
        \]
        In the nullary case the result follows from the identity laws in
        \(\mathbb B\).

        \item \textbf{Compare residual behaviours.}
        For a residual object
        \[
            r:u\to v,
        \]
        one has
        \[
        \begin{aligned}
            q_\theta(\delta_\theta(r,(x_i)_i))
            &=
            \theta_0(q_i(\delta_i(r,x_i)))_i,
        \end{aligned}
        \]
        which is the object value of
        \[
            \theta_{\mathbf{Beh}}(\beta_{P_i}(x_i))_i.
        \]
        The same calculation on residual arrows gives
        \[
            \omega_\theta(a,\delta_\theta(r,(x_i)_i))
            =
            \theta_1(\omega_i(a,\delta_i(r,x_i)))_i.
        \]
        Therefore
        \[
            \beta_{\theta(P_i)_i}
            =
            \theta_{\mathbf{Beh}}(\beta_{P_i})_i.
        \]
    \end{enumerate}
\end{proof}

\begin{lemma}[Structured operations preserve derivatives and canonical output]
\label{lem:operations-preserve-derivatives-output}
    For every operation
    \[
        \theta:\mathbb B^n\to\mathbb B
    \]
    and every arrow
    \[
        a:u\to v
    \]
    of \(\mathbb A\),
    \[
        \partial_a\bigl(\theta_{\mathbf{Beh}}(H_i)_i\bigr)
        =
        \theta_{\mathbf{Beh}}(\partial_aH_i)_i
    \]
    and
    \[
        \omega_{\mathbf{Beh}}
        \bigl(a,\theta_{\mathbf{Beh}}(H_i)_i\bigr)
        =
        \theta_1\bigl(\omega_{\mathbf{Beh}}(a,H_i)\bigr)_i.
    \]
\end{lemma}

\begin{proof}
    \leavevmode
    \begin{enumerate}
        \item \textbf{Use precomposition for derivatives.}
        Since
        \[
            \partial_aH=H\circ a_*,
        \]
        one has
        \[
        \begin{aligned}
            \partial_a\bigl(\theta\circ\langle H_i\rangle_i\bigr)
            &=
            \bigl(\theta\circ\langle H_i\rangle_i\bigr)\circ a_* \\
            &=
            \theta\circ\langle H_i\circ a_*\rangle_i \\
            &=
            \theta_{\mathbf{Beh}}(\partial_aH_i)_i.
        \end{aligned}
        \]
        The commutation is
        \[
        \begin{tikzcd}[column sep=huge, row sep=large]
            \mathbb A/u
                \arrow[r, "a_*"]
                \arrow[dr, "{\langle \partial_aH_i\rangle_i}"']
            &
            \mathbb A/v
                \arrow[d, "{\langle H_i\rangle_i}"]
            \\
            &
            \mathbb B^n
                \arrow[d, "\theta"]
            \\
            &
            \mathbb B .
        \end{tikzcd}
        \]

        \item \textbf{Use functoriality for canonical output.}
        The canonical output is the value on the residual morphism
        \[
            a\to 1_v.
        \]
        Applying
        \[
            \theta\circ\langle H_i\rangle_i
        \]
        gives
        \[
            \theta_1(H_i(a,1_v))_i
            =
            \theta_1(\omega_{\mathbf{Beh}}(a,H_i))_i.
        \]
        Diagrammatically,
        \[
        \begin{tikzcd}[column sep=large, row sep=large]
            (H_i(a))_i
                \arrow[r, "{(H_i(a\to 1_v))_i}"]
                \arrow[d, "\theta_0"']
            &
            (H_i(1_v))_i
                \arrow[d, "\theta_0"]
            \\
            \theta(H_i(a))_i
                \arrow[r, "{\theta_1(H_i(a\to 1_v))_i}"']
            &
            \theta(H_i(1_v))_i .
        \end{tikzcd}
        \]

        \item \textbf{Handle nullary operations.}
        In the nullary case both claims say that a constant functor remains
        constant under residual reindexing and sends residual morphisms to the
        identity arrow selected by the nullary functor.
    \end{enumerate}
\end{proof}

\subsection{\texorpdfstring{\(\mathbb T\)}{T}-coequational theories}
\label{subsec:T-coequational-theories}

\begin{definition}[\(\mathbb T\)-coequational theory]
\label{def:T-coequational-theory}
    Let \(\mathbb B\) be a \(\mathbb T\)-structured output category. A
    \textbf{local \(\mathbb T\)-coequational theory} is a local coequational
    theory
    \[
        V\hookrightarrow\mathbf{Beh}^{\mathrm{can}}_{\mathbb A,\mathbb B}
    \]
    closed under the pointwise operations induced by the \(\mathbb T\)-structure
    on \(\mathbb B\), where each operation is regarded with its induced
    \(Z\)-coordinate.

    Explicitly, for an \(n\)-ary operation
    \[
        \theta:\mathbb B^n\to\mathbb B,
    \]
    define, for \(n>0\),
    \[
        V^{(n)}_\theta
        :=
        \underbrace{V\times_{A_0}\cdots\times_{A_0}V}_{n},
    \]
    equipped with the \(\theta\)-induced map to
    \[
        Z=A_0\times B_0
    \]
    given by
    \[
        (H_1,\dots,H_n)
        \longmapsto
        \Bigl(
            p_V(H_1),
            \theta_0(q_V(H_1),\dots,q_V(H_n))
        \Bigr).
    \]
    Equivalently, \(V^{(n)}_\theta\) is the subobject of
    \[
        \mathbf{Beh}_0^{(n)}
        =
        \mathbf{Beh}_0\times_{A_0}\cdots\times_{A_0}\mathbf{Beh}_0
    \]
    whose components lie in \(V\), regarded over \(Z\) by the displayed
    \(\theta\)-induced output coordinate.

    For \(n=0\), set
    \[
        V^{(0)}_\theta:=A_0
    \]
    with \(Z\)-coordinate
    \[
        v\longmapsto (v,b_\theta).
    \]
    Closure under \(\theta\) means that
    \[
        \theta_{\mathbf{Beh}}:V^{(n)}_\theta\to\mathbf{Beh}_0
    \]
    factors through
    \[
        j_V:V\hookrightarrow\mathbf{Beh}_0
    \]
    as a map over \(Z\). The induced map is denoted
    \[
        \theta_V:V^{(n)}_\theta\to V.
    \]
\end{definition}

The closure condition is displayed by
\[
\begin{tikzcd}[column sep=large, row sep=large]
    V^{(n)}_\theta
        \arrow[rr, "{\theta_{\mathbf{Beh}}}"]
        \arrow[dr, dashed, "\theta_V"']
    &&
    \mathbf{Beh}_0
    \\
    &
    V
        \arrow[ur, hook, "j_V"']
    &
\end{tikzcd}
\qquad
\text{in }\mathcal E/Z.
\]
The compatibility between algebraic closure and derivative closure is
summarized by the cube-like diagram
\[
\begin{tikzcd}[column sep=huge, row sep=large]
    A_1\times_{A_0}V^{(n)}_\theta
        \arrow[r, "{1\times \theta_V}"]
        \arrow[d, "\delta_{V^{(n)}_\theta}"']
        \arrow[dr, phantom, "\circlearrowleft", very near start]
    &
    A_1\times_{A_0}V
        \arrow[d, "\delta_V"]
    \\
    V^{(n)}_\theta
        \arrow[r, "\theta_V"']
        \arrow[d, hook]
    &
    V
        \arrow[d, hook, "j_V"]
    \\
    \mathbf{Beh}^{(n)}_{\theta,0}
        \arrow[r, "\theta_{\mathbf{Beh}}"']
    &
    \mathbf{Beh}_0 .
\end{tikzcd}
\]

\subsection{Semantic images and structured operations}
\label{subsec:semantic-images-structured-operations}

\begin{lemma}[Semantic images under structured operations]
\label{lem:semantic-images-structured-operations}
    Let
    \[
        \theta:\mathbb B^n\to\mathbb B
    \]
    be an operation of \(\mathbb T\), and let
    \[
        P_1,\dots,P_n
    \]
    be Mealy morphisms. For \(n>0\), the semantic quotient maps
    \[
        P_i\twoheadrightarrow\operatorname{Beh}(P_i)
    \]
    induce a regular epimorphism
    \[
        \theta(P_1,\dots,P_n)
        \twoheadrightarrow
        \operatorname{Beh}(P_1)\times_{A_0}\cdots\times_{A_0}
        \operatorname{Beh}(P_n),
    \]
    where the codomain is equipped with the \(\theta\)-induced
    \(Z\)-coordinate. Moreover,
    \[
        \operatorname{Im}
        \left(
        \theta_{\mathbf{Beh}}:
        \operatorname{Beh}(P_1)\times_{A_0}\cdots\times_{A_0}
        \operatorname{Beh}(P_n)
        \to
        \mathbf{Beh}_0
        \right)
        =
        \operatorname{Beh}(\theta(P_1,\dots,P_n)).
    \]
    For \(n=0\), the image of
    \[
        A_0\to\mathbf{Beh}_0
    \]
    induced by the constant residual behaviour is
    \[
        \operatorname{Beh}(\theta()).
    \]
\end{lemma}

\begin{proof}
    \leavevmode
    \begin{enumerate}
        \item \textbf{Form the product of semantic quotients.}
        For \(n>0\), the product of the regular epimorphisms
        \[
            P_i\twoheadrightarrow\operatorname{Beh}(P_i)
        \]
        over \(A_0\) is a regular epimorphism in the topos slice. It gives
        \[
            P_1\times_{A_0}\cdots\times_{A_0}P_n
            \twoheadrightarrow
            \operatorname{Beh}(P_1)\times_{A_0}\cdots\times_{A_0}
            \operatorname{Beh}(P_n).
        \]
        The source is the carrier of \(\theta(P_1,\dots,P_n)\), and the target
        has the \(\theta\)-induced \(Z\)-coordinate. The construction is
        \[
        \begin{tikzcd}[column sep=large, row sep=large]
            P_1\times_{A_0}\cdots\times_{A_0}P_n
                \arrow[r, two heads]
                \arrow[d]
            &
            \operatorname{Beh}(P_1)\times_{A_0}\cdots\times_{A_0}
            \operatorname{Beh}(P_n)
                \arrow[d]
            \\
            A_0
                \arrow[r, equal]
            &
            A_0 .
        \end{tikzcd}
        \]

        \item \textbf{Compare the semantic composites.}
        By Lemma~\ref{lem:structured-operations-semantic},
        \[
            \beta_{\theta(P_i)_i}
            =
            \theta_{\mathbf{Beh}}(\beta_{P_i})_i.
        \]
        Hence the semantic map of \(\theta(P_1,\dots,P_n)\) is the composite of
        the product of semantic quotient maps with
        \[
            \theta_{\mathbf{Beh}}:
            \operatorname{Beh}(P_1)\times_{A_0}\cdots\times_{A_0}
            \operatorname{Beh}(P_n)
            \to
            \mathbf{Beh}_0.
        \]
        This gives
        \[
        \begin{tikzcd}[column sep=large, row sep=large]
            \theta(P_i)_i
                \arrow[r, two heads]
                \arrow[rr, bend left=18, "\beta_{\theta(P_i)_i}"]
            &
            \prod_{A_0}\operatorname{Beh}(P_i)
                \arrow[r, "\theta_{\mathbf{Beh}}"]
            &
            \mathbf{Beh}_0 .
        \end{tikzcd}
        \]

        \item \textbf{Compare regular images.}
        Precomposition with a regular epimorphism does not change the regular
        image. Therefore the image of the displayed \(\theta_{\mathbf{Beh}}\)
        is exactly the semantic image of
        \[
            \theta(P_1,\dots,P_n).
        \]

        \item \textbf{Handle nullary operations.}
        For \(n=0\), the semantic map of \(\theta()\) is the map
        \[
            A_0\to\mathbf{Beh}_0
        \]
        assigning to \(v\) the constant residual behaviour
        \[
            \mathbb A/v\to 1\xrightarrow{\theta}\mathbb B.
        \]
        Its regular image is, by definition,
        \[
            \operatorname{Beh}(\theta()).
        \]
    \end{enumerate}
\end{proof}

\subsection{\texorpdfstring{\(\mathbb T\)}{T}-coequational closure}
\label{subsec:T-coequational-closure}

\begin{theorem}[\(\mathbb T\)-coequational closure]
\label{thm:T-coequational-closure}
    Let
    \[
        V\hookrightarrow\mathbf{Beh}^{\mathrm{can}}
    \]
    be a local \(\mathbb T\)-coequational theory. Then
    \[
        \mathsf{Mod}(V)
    \]
    is closed under strict homomorphic preimages, regular-epimorphic quotients,
    small coproducts, and the operations induced by \(\mathbb T\).
\end{theorem}

\begin{proof}
    \leavevmode
    \begin{enumerate}
        \item \textbf{Use local closure.}
        Closure under strict homomorphic preimages, regular-epimorphic
        quotients, and small coproducts is
        Theorem~\ref{thm:local-coequational-closure}.

        \item \textbf{Use semantic compatibility.}
        Let
        \[
            P_1,\dots,P_n\in\mathsf{Mod}(V).
        \]
        Choose strict factorizations
        \[
            \bar\beta_{P_i}:P_i\to V
        \]
        with
        \[
            j_V\bar\beta_{P_i}=\beta_{P_i}.
        \]
        These induce a map from the carrier of
        \[
            \theta(P_1,\dots,P_n)
        \]
        to \(V^{(n)}_\theta\):
        \[
        \begin{tikzcd}[column sep=large, row sep=large]
            P_1\times_{A_0}\cdots\times_{A_0}P_n
                \arrow[r, "{(\bar\beta_{P_i})_i}"]
                \arrow[dr, "{(\beta_{P_i})_i}"']
            &
            V^{(n)}_\theta
                \arrow[d, hook]
            \\
            &
            \mathbf{Beh}^{(n)}_{\theta,0}.
        \end{tikzcd}
        \]

        \item \textbf{Use algebraic closure of \(V\).}
        Since \(V\) is closed under \(\theta\), there is a map
        \[
            \theta_V:V^{(n)}_\theta\to V
        \]
        with
        \[
            j_V\theta_V
            =
            \theta_{\mathbf{Beh}}\big|_{V^{(n)}_\theta}.
        \]
        By Lemma~\ref{lem:structured-operations-semantic},
        \[
            \beta_{\theta(P_i)_i}
            =
            \theta_{\mathbf{Beh}}(\beta_{P_i})_i.
        \]
        Hence
        \[
            \beta_{\theta(P_i)_i}
        \]
        factors through \(V\), so
        \[
            \theta(P_1,\dots,P_n)\in\mathsf{Mod}(V).
        \]
        The factorization is
        \[
        \begin{tikzcd}[column sep=large, row sep=large]
            \theta(P_i)_i
                \arrow[r, "{(\bar\beta_{P_i})_i}"]
                \arrow[rrr, bend left=18, "\beta_{\theta(P_i)_i}"]
            &
            V^{(n)}_\theta
                \arrow[r, "\theta_V"]
            &
            V
                \arrow[r, hook, "j_V"]
            &
            \mathbf{Beh}^{\mathrm{can}} .
        \end{tikzcd}
        \]

        \item \textbf{Handle nullary operations.}
        The nullary case is the same argument using
        \[
            V^{(0)}_\theta=A_0
        \]
        and the closure map
        \[
            A_0\to V.
        \]
    \end{enumerate}
\end{proof}

\subsection{The structured Birkhoff theorem}
\label{subsec:T-coequational-birkhoff}

\begin{theorem}[\(\mathbb T\)-coequational Birkhoff theorem]
\label{thm:T-coequational-birkhoff}
    Let \(\mathbb B\) be a \(\mathbb T\)-structured output category, and let
    \[
        \mathcal C\subseteq\mathsf{Mly}^{\mathrm{str}}_{\mathbb A,\mathbb B}
    \]
    be a replete full subcategory. Then \(\mathcal C\) is of the form
    \[
        \mathcal C=\mathsf{Mod}(V)
    \]
    for a local \(\mathbb T\)-coequational theory
    \[
        V\hookrightarrow\mathbf{Beh}^{\mathrm{can}}_{\mathbb A,\mathbb B}
    \]
    if and only if \(\mathcal C\) is closed under:
    \begin{enumerate}
        \item strict homomorphic preimages;
        \item regular-epimorphic quotients;
        \item small coproducts;
        \item the operations induced by \(\mathbb T\).
    \end{enumerate}
\end{theorem}

\begin{proof}
    \leavevmode
    \begin{enumerate}
        \item \textbf{Prove necessity.}
        If
        \[
            \mathcal C=\mathsf{Mod}(V)
        \]
        for a local \(\mathbb T\)-coequational theory \(V\), then the closure
        properties are exactly Theorem~\ref{thm:T-coequational-closure}.

        \item \textbf{Construct the local candidate.}
        Conversely, assume \(\mathcal C\) has the four closure properties. By
        Theorem~\ref{thm:local-coequational-birkhoff},
        \[
            V_{\mathcal C}
            =
            \bigvee_{P\in\mathcal C}\operatorname{Beh}(P)
        \]
        is a local coequational theory and
        \[
            \mathcal C=\mathsf{Mod}(V_{\mathcal C}).
        \]
        It remains to prove that \(V_{\mathcal C}\) is closed under the
        pointwise \(\mathbb T\)-operations.

        Choose a small representative set
        \[
            \mathcal S_{\mathcal C}
            \subseteq
            \operatorname{Sub}_{\mathcal E/Z}(\mathbf{Beh}_0)
        \]
        for the semantic-image subobjects
        \[
            \operatorname{Beh}(P),
            \qquad
            P\in\mathcal C.
        \]
        For each
        \[
            S\in\mathcal S_{\mathcal C},
        \]
        choose
        \[
            P_S\in\mathcal C
        \]
        with
        \[
            S=\operatorname{Beh}(P_S).
        \]

        \item \textbf{Cover the operation domain.}
        For an \(n\)-ary operation \(\theta\), the object
        \[
            (V_{\mathcal C})^{(n)}_\theta
        \]
        is formed over \(A_0\), with \(Z\)-coordinate induced by \(\theta_0\).
        Since
        \[
            V_{\mathcal C}
            =
            \bigvee_{S\in\mathcal S_{\mathcal C}}S,
        \]
        and since pullback preserves joins while finite meets distribute over
        small joins in a topos slice, the subobject of tuples whose components
        lie in \(V_{\mathcal C}\) is the join of the subobjects
        \[
            S_1\times_{A_0}\cdots\times_{A_0}S_n,
            \qquad
            S_i\in\mathcal S_{\mathcal C}.
        \]
        Therefore
        \[
            (V_{\mathcal C})^{(n)}_\theta
        \]
        is covered by the coproduct of these products, equipped with the same
        \(\theta\)-induced \(Z\)-coordinate:
        \[
        \begin{tikzcd}[column sep=large, row sep=large]
            \displaystyle
            \coprod_{(S_1,\dots,S_n)}
            S_1\times_{A_0}\cdots\times_{A_0}S_n
                \arrow[r, two heads]
                \arrow[dr]
            &
            (V_{\mathcal C})^{(n)}_\theta
                \arrow[d, hook]
            \\
            &
            \mathbf{Beh}^{(n)}_{\theta,0}.
        \end{tikzcd}
        \]

        \item \textbf{Compute on each cover component.}
        For each tuple
        \[
            (S_1,\dots,S_n)
            \in
            \mathcal S_{\mathcal C}^n,
        \]
        Lemma~\ref{lem:semantic-images-structured-operations} gives
        \[
            \theta_{\mathbf{Beh}}
            \bigl(
            S_1\times_{A_0}\cdots\times_{A_0}S_n
            \bigr)
            \leq
            \operatorname{Beh}
            \bigl(
            \theta(P_{S_1},\dots,P_{S_n})
            \bigr).
        \]
        Since \(\mathcal C\) is closed under the operations induced by
        \(\mathbb T\),
        \[
            \theta(P_{S_1},\dots,P_{S_n})\in\mathcal C,
        \]
        and hence
        \[
            \operatorname{Beh}
            \bigl(
            \theta(P_{S_1},\dots,P_{S_n})
            \bigr)
            \leq
            V_{\mathcal C}.
        \]
        The relevant factorization on each piece is
\[
\begin{tikzcd}[column sep=large, row sep=large]
    S_1\times_{A_0}\cdots\times_{A_0}S_n
        \arrow[r, "{\theta_{\mathbf{Beh}}}"]
        \arrow[dr, dashed, "{\theta_{V_{\mathcal C}}|_{S_1,\dots,S_n}}"']
    &
    \mathbf{Beh}_0
    \\
    &
    V_{\mathcal C}
        \arrow[u, hook, "j_{V_{\mathcal C}}"']
\end{tikzcd}
\]

        \item \textbf{Descend the factorization.}
        The operation map
        \[
            \theta_{\mathbf{Beh}}:
            (V_{\mathcal C})^{(n)}_\theta\to\mathbf{Beh}_0
        \]
        factors through \(V_{\mathcal C}\) after precomposition with a
        regular-epimorphic cover. By
        Lemma~\ref{lem:regular-epi-descent-for-subobjects-membership}, it factors through
        \(V_{\mathcal C}\). The descent diagram is
        \[
        \begin{tikzcd}[column sep=large, row sep=large]
            \displaystyle
            \coprod_{(S_1,\dots,S_n)}
            S_1\times_{A_0}\cdots\times_{A_0}S_n
                \arrow[r, two heads]
                \arrow[dr]
                \arrow[ddr, bend right=18, "\theta_{\mathbf{Beh}}"']
            &
            (V_{\mathcal C})^{(n)}_\theta
                \arrow[d, dashed]
                \arrow[dd, bend left=18, "\theta_{\mathbf{Beh}}"]
            \\
            &
            V_{\mathcal C}
                \arrow[d, hook]
            \\
            &
            \mathbf{Beh}_0 .
        \end{tikzcd}
        \]

        \item \textbf{Handle nullary operations.}
        For \(n=0\), closure of \(\mathcal C\) under the nullary operation
        gives
        \[
            \theta()\in\mathcal C.
        \]
        Hence
        \[
            \operatorname{Beh}(\theta())\leq V_{\mathcal C}.
        \]
        Since the constant behaviour map
        \[
            A_0\to\mathbf{Beh}_0
        \]
        factors through \(\operatorname{Beh}(\theta())\), it factors through
        \(V_{\mathcal C}\). Therefore the nullary operation also restricts to
        \(V_{\mathcal C}\).

        \item \textbf{Conclude.}
        Thus \(V_{\mathcal C}\) is a local \(\mathbb T\)-coequational theory,
        and
        \[
            \mathcal C=\mathsf{Mod}(V_{\mathcal C}).
        \]
    \end{enumerate}
\end{proof}

\subsection{Running specializations}
\label{subsec:structured-output-running-specializations}

For Boolean output, take
\[
    \mathbb B=2_{\mathrm{ch}},
\]
the chaotic internal category on a Boolean algebra object \(2\). The Boolean
operations
\[
    \wedge,\vee,\neg,0,1
\]
are internal functors
\[
    2_{\mathrm{ch}}^n\to 2_{\mathrm{ch}}.
\]
Since \(2_{\mathrm{ch}}\) is chaotic, output-comparison arrows are determined
by their endpoints.

In \(\mathbf{Set}\), a local Boolean coequational theory over \(\Sigma\) is a
Boolean subalgebra of languages over \(\Sigma^*\) closed under residual
derivatives:
\[
    L\longmapsto \partial_uL.
\]
Regularity is an additional finite-state condition, or an explicit restriction
to regular languages. Thus the derivative-closed Boolean structure is separated
from finite recognizability, as in the classical distinction between all
languages and regular languages
\cite{Eilenberg1974,HopcroftMotwaniUllman2006,Brzozowski1964}.

In the stateless specialization, if
\[
    F_1,\dots,F_n:\mathbb A\to\mathbb B
\]
are internal functors and
\[
    \theta:\mathbb B^n\to\mathbb B
\]
is a structural operation, then the induced stateless recognizer is the
internal functor
\[
    \mathbb A
    \xrightarrow{\langle F_i\rangle_i}
    \mathbb B^n
    \xrightarrow{\theta}
    \mathbb B.
\]
The residual slice functors of this composite are obtained by applying
\(\theta\) pointwise to the residual slice functors of the \(F_i\).

\section{Global input reindexing}
\label{sec:global-input-reindexing}

\subsection{Reindexing residual behaviours along input functors}
\label{subsec:input-reindexing-behaviours}

This section explains how local coequational theories vary with the input
category. This is the categorical source of closure under inverse morphic
preimage in the classical one-object setting, a closure operation appearing in
Eilenberg-type correspondences for varieties of languages
\cite{Eilenberg1974,AdamekMiliusMyersUrbat2019TOCL}.

Fix \(\mathbb B\), and, when desired, a finite-product theory \(\mathbb T\) for
which \(\mathbb B\) is a \(\mathbb T\)-structured output category. Let
\[
    \mathsf{Inp}
\]
be a chosen category of internal input categories and internal functors. For
every
\[
    \mathbb A\in\mathsf{Inp},
\]
there is a canonical behaviour object
\[
    \mathbf{Beh}_{\mathbb A,\mathbb B,0}
\]
and a strict Mealy category
\[
    \mathsf{Mly}^{\mathrm{str}}_{\mathbb A,\mathbb B}.
\]
We write
\[
    Z_{\mathbb A}:=A_0\times B_0.
\]

Let
\[
    F:\mathbb A'\to\mathbb A
\]
be an internal functor. For each \(v'\in A'_0\), there is an induced residual
slice functor
\[
    F_{/v'}:
    \mathbb A'/v'\to \mathbb A/F_0(v')
\]
sending
\[
    r':u'\to v'
\]
to
\[
    F_1(r'):F_0(u')\to F_0(v').
\]
On residual triangles it acts by applying \(F\):
\[
\begin{tikzcd}[column sep=large, row sep=large]
    w'
        \arrow[r, "a'"]
        \arrow[rr, "r'\circ a'", bend left=15]
    &
    u'
        \arrow[r, "r'"']
    &
    v'
\end{tikzcd}
\qquad\longmapsto\qquad
\begin{tikzcd}[column sep=large, row sep=large]
    F_0w'
        \arrow[r, "F_1a'"]
        \arrow[rr, "F_1(r'\circ a')"', bend right=15]
    &
    F_0u'
        \arrow[r, "F_1r'"]
    &
    F_0v'.
\end{tikzcd}
\]
Functoriality of \(F\) gives
\[
    F_1(r'\circ a')=F_1(r')\circ F_1(a').
\]

By the representing property of the behaviour object, these residual slice
functors induce a reindexing map
\[
    F^\sharp:
    A'_0\times_{F_0,A_0,p_{\mathbf{Beh}}}
    \mathbf{Beh}_{\mathbb A,\mathbb B,0}
    \to
    \mathbf{Beh}_{\mathbb A',\mathbb B,0}
\]
given by
\[
    F^\sharp(v',H)=H\circ F_{/v'}.
\]
The domain is regarded over
\[
    Z_{\mathbb A'}=A'_0\times B_0
\]
by
\[
    (v',H)\longmapsto (v',q_{\mathbf{Beh}}(H)).
\]
With this convention, \(F^\sharp\) is a map over \(Z_{\mathbb A'}\):
\[
\begin{tikzcd}[column sep=large, row sep=large]
    A'_0\times_{A_0}\mathbf{Beh}_{\mathbb A,\mathbb B,0}
        \arrow[r, "F^\sharp"]
        \arrow[d, "{(v',H)\mapsto(v',qH)}"']
    &
    \mathbf{Beh}_{\mathbb A',\mathbb B,0}
        \arrow[d, "{(p,q)}"]
    \\
    A'_0\times B_0
        \arrow[r, equal]
    &
    A'_0\times B_0 .
\end{tikzcd}
\]

\begin{lemma}[Reindexing preserves canonical structure]
\label{lem:global-reindexing-preserves-structure}
    The map \(F^\sharp\) preserves derivatives and canonical output
    comparisons. Explicitly,
    \[
        \partial_{a'}(F^\sharp(v',H))
        =
        F^\sharp(u',\partial_{F_1(a')}H)
    \]
    for every
    \[
        a':u'\to v',
    \]
    and
    \[
        \omega_{\mathbf{Beh}_{\mathbb A',\mathbb B}}
        (a',F^\sharp(v',H))
        =
        \omega_{\mathbf{Beh}_{\mathbb A,\mathbb B}}
        (F_1(a'),H).
    \]
    If \(\mathbb B\) is \(\mathbb T\)-structured, then \(F^\sharp\) also
    preserves the pointwise \(\mathbb T\)-structure.
\end{lemma}

\begin{proof}
    \leavevmode
    \begin{enumerate}
        \item \textbf{Compare residual transports.}
        The residual slice functors commute with postcomposition:
        \[
            F_{/v'}\circ a'_*
            =
            (F_1a')_*\circ F_{/u'}.
        \]
        Diagrammatically,
        \[
        \begin{tikzcd}[column sep=huge, row sep=large]
            \mathbb A'/u'
                \arrow[r, "a'_*"]
                \arrow[d, "F_{/u'}"']
            &
            \mathbb A'/v'
                \arrow[d, "F_{/v'}"]
            \\
            \mathbb A/F_0u'
                \arrow[r, "{(F_1a')_*}"']
            &
            \mathbb A/F_0v'.
        \end{tikzcd}
        \]
        Precomposition with this square gives the derivative identity:
        \[
        \begin{tikzcd}[column sep=huge, row sep=large]
            \mathbb A'/u'
                \arrow[r, "a'_*"]
                \arrow[d, "F_{/u'}"']
                \arrow[dr, phantom, "\circlearrowleft", very near start]
            &
            \mathbb A'/v'
                \arrow[d, "F_{/v'}"]
                \arrow[ddr, "{H\circ F_{/v'}}", bend left=18]
            \\
            \mathbb A/F_0u'
                \arrow[r, "{(F_1a')_*}"']
                \arrow[drr, "{\partial_{F_1a'}H}"', bend right=18]
            &
            \mathbb A/F_0v'
                \arrow[dr, "H"]
            \\
            &&
            \mathbb B .
        \end{tikzcd}
        \]

        \item \textbf{Compare canonical outputs.}
        The canonical output of \(F^\sharp(v',H)\) along \(a'\) is the arrow
        assigned by
        \[
            H\circ F_{/v'}
        \]
        to the residual morphism
        \[
            a'\to 1_{v'}
        \]
        in \(\mathbb A'/v'\). Under \(F_{/v'}\), this residual morphism is sent
        to
        \[
            F_1(a')\to 1_{F_0(v')}
        \]
        in \(\mathbb A/F_0(v')\). Hence the output equality follows:
        \[
        \begin{tikzcd}[column sep=large, row sep=large]
            a'
                \arrow[r]
                \arrow[d, "{F_{/v'}}"']
            &
            1_{v'}
                \arrow[d, "{F_{/v'}}"]
            \\
            F_1(a')
                \arrow[r]
            &
            1_{F_0(v')} .
        \end{tikzcd}
        \]

        \item \textbf{Compare pointwise operations.}
        The map \(F^\sharp\) acts by precomposition, while the pointwise
        \(\mathbb T\)-structure is defined by postcomposition with the
        structural functors
        \[
            \mathbb B^n\to\mathbb B.
        \]
        Therefore the two operations commute. The diagram is
        \[
        \begin{tikzcd}[column sep=huge, row sep=large]
            \mathbb A'/v'
                \arrow[r, "F_{/v'}"]
                \arrow[dr, "{\langle H_i\circ F_{/v'}\rangle_i}"']
            &
            \mathbb A/F_0v'
                \arrow[d, "{\langle H_i\rangle_i}"]
            \\
            &
            \mathbb B^n
                \arrow[d, "\theta"]
            \\
            &
            \mathbb B .
        \end{tikzcd}
        \]
    \end{enumerate}
\end{proof}

\subsection{Input pullback of Mealy morphisms}
\label{subsec:input-pullback-mealy}

\begin{definition}[Input pullback]
\label{def:input-pullback-mealy}
    Let
    \[
        P:\mathbb A\rightsquigarrow\mathbb B
    \]
    be a Mealy morphism. Define
    \[
        F^*P:\mathbb A'\rightsquigarrow\mathbb B
    \]
    as follows. The carrier is
    \[
        F^*P
        :=
        A'_0\times_{F_0,A_0,p_P}P.
    \]
    A state is a pair
    \[
        (v',x)
    \]
    with
    \[
        F_0(v')=p_P(x).
    \]
    The input projection is
    \[
        p_{F^*P}(v',x)=v',
    \]
    and the output projection is
    \[
        q_{F^*P}(v',x)=q_P(x).
    \]
    For an arrow
    \[
        a':u'\to v',
    \]
    define
    \[
        \delta_{F^*P}(a',(v',x))
        =
        (u',\delta_P(F_1a',x)),
    \]
    and
    \[
        \omega_{F^*P}(a',(v',x))
        =
        \omega_P(F_1a',x).
    \]
\end{definition}

The carrier is the pullback
\[
\begin{tikzcd}[column sep=large, row sep=large]
    F^*P
        \arrow[r, "\pi_P"]
        \arrow[d]
        \arrow[dr, phantom, "\lrcorner", very near start]
    &
    P
        \arrow[d, "p_P"]
    \\
    A'_0
        \arrow[r, "F_0"']
    &
    A_0 .
\end{tikzcd}
\]
The transition is induced by the square
\[
\begin{tikzcd}[column sep=huge, row sep=large]
    A'_1\times_{A'_0}F^*P
        \arrow[r, "{F_1\times \pi_P}"]
        \arrow[d, "\delta_{F^*P}"']
    &
    A_1\times_{A_0}P
        \arrow[d, "\delta_P"]
    \\
    F^*P
        \arrow[r, "\pi_P"']
    &
    P ,
\end{tikzcd}
\]
together with the input-object component \(u'\).

For a strict semantic homomorphism
\[
    f:P\to Q
\]
define
\[
    F^*f:F^*P\to F^*Q
\]
by
\[
    F^*f(v',x)=(v',fx).
\]
This assignment is functorial.

\begin{proposition}[Input pullback and semantics]
\label{prop:input-pullback-mealy}
    The data above define a Mealy morphism
    \[
        F^*P:\mathbb A'\rightsquigarrow\mathbb B.
    \]
    Moreover,
    \[
        F^*:\mathsf{Mly}^{\mathrm{str}}_{\mathbb A,\mathbb B}
        \to
        \mathsf{Mly}^{\mathrm{str}}_{\mathbb A',\mathbb B}
    \]
    is a functor, and the semantic map satisfies
    \[
        \beta_{F^*P}
        =
        F^\sharp
        \circ
        \langle p_{F^*P},\,\beta_P\pi_P\rangle.
    \]
    In generalized elements,
    \[
        \beta_{F^*P}(v',x)
        =
        F^\sharp(v',\beta_P(x)).
    \]
\end{proposition}

\begin{proof}
    \leavevmode
    \begin{enumerate}
        \item \textbf{Check typing.}
        If
        \[
            a':u'\to v',
        \]
        then
        \[
            F_1(a'):F_0(u')\to F_0(v').
        \]
        Hence \(\delta_P(F_1a',x)\) is defined whenever
        \(p_P(x)=F_0(v')\), and lies over \(F_0(u')\). Therefore
        \[
            (u',\delta_P(F_1a',x))
        \]
        lies in
        \[
            A'_0\times_{A_0}P.
        \]
        The typing square is
        \[
        \begin{tikzcd}[column sep=large, row sep=large]
            \delta_P(F_1a',x)
                \arrow[r, maps to]
                \arrow[d, maps to]
            &
            F_0(u')
                \arrow[d, "F_0"]
            \\
            x
                \arrow[r, maps to]
            &
            F_0(v') .
        \end{tikzcd}
        \]

        \item \textbf{Verify the Mealy laws.}
        The unit and multiplication laws follow from functoriality of \(F\) and
        the Mealy laws in \(P\). The output laws are inherited in the same way.
        For instance, for \(a':u'\to v'\) and \(b':v'\to w'\),
        \[
        \begin{aligned}
            \delta_{F^*P}(b'\circ a',(w',x))
            &=
            (u',\delta_P(F_1(b'\circ a'),x)) \\
            &=
            (u',\delta_P(F_1b'\circ F_1a',x)) \\
            &=
            (u',\delta_P(F_1a',\delta_P(F_1b',x))).
        \end{aligned}
        \]

        \item \textbf{Check the action on strict semantic homomorphisms.}
        Let
        \[
            f:P\to Q
        \]
        be strict. The map
        \[
            F^*f(v',x)=(v',fx)
        \]
        lies over \(A'_0\) and \(B_0\). Transition preservation follows from
        \[
        \begin{aligned}
            F^*f(\delta_{F^*P}(a',(v',x)))
            &=
            F^*f(u',\delta_P(F_1a',x)) \\
            &=
            (u',f\delta_P(F_1a',x)) \\
            &=
            (u',\delta_Q(F_1a',fx)) \\
            &=
            \delta_{F^*Q}(a',(v',fx)).
        \end{aligned}
        \]
        Output preservation follows from strictness of \(f\):
        \[
            \omega_P(F_1a',x)=\omega_Q(F_1a',fx).
        \]
        Identities and composition are inherited from \(\mathcal E\), so
        \(F^*\) is a functor. The strictness diagram is
        \[
        \begin{tikzcd}[column sep=huge, row sep=large]
            A'_1\times_{A'_0}F^*P
                \arrow[r, "{1\times F^*f}"]
                \arrow[d, "\delta_{F^*P}"']
            &
            A'_1\times_{A'_0}F^*Q
                \arrow[d, "\delta_{F^*Q}"]
            \\
            F^*P
                \arrow[r, "F^*f"']
            &
            F^*Q .
        \end{tikzcd}
        \]

        \item \textbf{Identify residual semantics.}
        Residual execution in \(F^*P\) along a residual arrow \(r'\) is
        execution in \(P\) along \(F_1(r')\). Thus the residual behaviour of
        \((v',x)\) is the residual behaviour of \(x\) precomposed with
        \[
            F_{/v'}:\mathbb A'/v'\to\mathbb A/F_0(v').
        \]
        Hence
        \[
            \beta_{F^*P}(v',x)
            =
            F^\sharp(v',\beta_P(x)).
        \]

        \item \textbf{Display the semantic triangle.}
        This identity is represented by
        \[
        \begin{tikzcd}[column sep=large, row sep=large]
            F^*P
                \arrow[r, "{\langle p_{F^*P},\,\beta_P\pi_P\rangle}"]
                \arrow[dr, "\beta_{F^*P}"']
            &
            A'_0\times_{A_0}\mathbf{Beh}_{\mathbb A,\mathbb B,0}
                \arrow[d, "F^\sharp"]
            \\
            &
            \mathbf{Beh}_{\mathbb A',\mathbb B,0}.
        \end{tikzcd}
        \]
        With the carrier projection included, the same statement is
        \[
        \begin{tikzcd}[column sep=large, row sep=large]
            F^*P
                \arrow[r, "\pi_P"]
                \arrow[d, "\beta_{F^*P}"']
                \arrow[dr, phantom, "\circlearrowleft", very near start]
            &
            P
                \arrow[d, "\beta_P"]
            \\
            \mathbf{Beh}_{\mathbb A',\mathbb B,0}
            &
            \mathbf{Beh}_{\mathbb A,\mathbb B,0}
                \arrow[l, dashed, "F^\sharp"']
        \end{tikzcd}
        \]
        where \(F^\sharp\) is read on the pulled-back domain
        \(A'_0\times_{A_0}\mathbf{Beh}_{\mathbb A,\mathbb B,0}\).
    \end{enumerate}
\end{proof}

\begin{lemma}[Input pullback and semantic images]
\label{lem:input-pullback-semantic-images}
    Let
    \[
        F:\mathbb A'\to\mathbb A
    \]
    be an internal functor and let
    \[
        P:\mathbb A\rightsquigarrow\mathbb B.
    \]
    Then
    \[
        \operatorname{Im}
        \left(
            A'_0\times_{A_0}\operatorname{Beh}(P)
            \xrightarrow{F^\sharp}
            \mathbf{Beh}_{\mathbb A',\mathbb B,0}
        \right)
        =
        \operatorname{Beh}(F^*P).
    \]
\end{lemma}

\begin{proof}
    \leavevmode
    \begin{enumerate}
        \item \textbf{Pull back the semantic quotient.}
        Pull back the regular epimorphism
        \[
            P\twoheadrightarrow\operatorname{Beh}(P)
        \]
        along
        \[
            F_0:A'_0\to A_0.
        \]
        This gives a regular epimorphism
        \[
            A'_0\times_{A_0}P
            \twoheadrightarrow
            A'_0\times_{A_0}\operatorname{Beh}(P).
        \]
        The source is the carrier of \(F^*P\). The pullback is
        \[
        \begin{tikzcd}[column sep=large, row sep=large]
            A'_0\times_{A_0}P
                \arrow[r, two heads]
                \arrow[d]
                \arrow[dr, phantom, "\lrcorner", very near start]
            &
            A'_0\times_{A_0}\operatorname{Beh}(P)
                \arrow[d]
            \\
            P
                \arrow[r, two heads]
            &
            \operatorname{Beh}(P).
        \end{tikzcd}
        \]

        \item \textbf{Compare semantic maps.}
        By Proposition~\ref{prop:input-pullback-mealy}, the composite
        \[
            F^*P
            \to
            A'_0\times_{A_0}\operatorname{Beh}(P)
            \xrightarrow{F^\sharp}
            \mathbf{Beh}_{\mathbb A',\mathbb B,0}
        \]
        is
        \[
            \beta_{F^*P}.
        \]
        Thus
        \[
        \begin{tikzcd}[column sep=large, row sep=large]
            F^*P
                \arrow[r, two heads]
                \arrow[rr, bend left=18, "\beta_{F^*P}"]
            &
            A'_0\times_{A_0}\operatorname{Beh}(P)
                \arrow[r, "F^\sharp"]
            &
            \mathbf{Beh}_{\mathbb A',\mathbb B,0}
        \end{tikzcd}
        \]
        is the semantic factorization before taking the final regular image.

        \item \textbf{Pass to images.}
        Since the first map is a regular epimorphism, the image of the second
        map is the same as the image of the composite. The image of the
        composite is
        \[
            \operatorname{Beh}(F^*P).
        \]
        Therefore
        \[
            \operatorname{Im}
            \left(
                F^\sharp:
                A'_0\times_{A_0}\operatorname{Beh}(P)
                \to
                \mathbf{Beh}_{\mathbb A',\mathbb B,0}
            \right)
            =
            \operatorname{Beh}(F^*P).
        \]
    \end{enumerate}
\end{proof}

\subsection{Running specializations}
\label{subsec:global-reindexing-running-specializations}

In the one-object free-monoid case, an input functor
\[
    F:\mathbb A_{\Gamma}\to\mathbb A_{\Sigma}
\]
is a monoid homomorphism
\[
    h:\Gamma^*\to\Sigma^*.
\]
For a language
\[
    L:\Sigma^*\to 2,
\]
the reindexed behaviour is
\[
    F^\sharp(L)=L\circ h.
\]
Thus \(F^\sharp\) is inverse morphic preimage:
\[
    h^{-1}L(w)=L(h(w)).
\]

In the stateless case, input pullback sends an internal functor
\[
    G:\mathbb A\to\mathbb B
\]
to the composite
\[
    \mathbb A'
    \xrightarrow{F}
    \mathbb A
    \xrightarrow{G}
    \mathbb B.
\]
The residual slice functor of \(G\circ F\) at \(v'\) is obtained by
precomposing the residual slice functor of \(G\) at \(F_0(v')\) with
\[
    F_{/v'}:\mathbb A'/v'\to\mathbb A/F_0(v').
\]

\section{Global coequational theories}
\label{sec:global-coequational-theories}

\subsection{Global theories and global models}
\label{subsec:global-theories-and-models}

\begin{definition}[Global coequational theory]
\label{def:global-mealy-coequational-theory}
    A \textbf{global coequational theory} consists of a local coequational
    theory
    \[
        V_{\mathbb A}\hookrightarrow
        \mathbf{Beh}^{\mathrm{can}}_{\mathbb A,\mathbb B}
    \]
    for each \(\mathbb A\in\mathsf{Inp}\), such that for every internal functor
    \[
        F:\mathbb A'\to\mathbb A
    \]
    the reindexing map
    \[
        F^\sharp:
        A'_0\times_{F_0,A_0,p_{\mathbf{Beh}}}
        \mathbf{Beh}_{\mathbb A,\mathbb B,0}
        \to
        \mathbf{Beh}_{\mathbb A',\mathbb B,0}
    \]
    sends
    \[
        A'_0\times_{F_0,A_0,p_{V_{\mathbb A}}}V_{\mathbb A}
    \]
    into
    \[
        V_{\mathbb A'}.
    \]
    If \(\mathbb B\) is \(\mathbb T\)-structured, a \textbf{global
    \(\mathbb T\)-coequational theory} is such a family for which each
    \(V_{\mathbb A}\) is a local \(\mathbb T\)-coequational theory.
\end{definition}

The global reindexing condition is displayed by
\[
\begin{tikzcd}[column sep=large, row sep=large]
    A'_0\times_{A_0}V_{\mathbb A}
        \arrow[rr, "{F^\sharp|_{V_{\mathbb A}}}"]
        \arrow[dr, dashed]
    &&
    \mathbf{Beh}_{\mathbb A',\mathbb B,0}
    \\
    &
    V_{\mathbb A'}
        \arrow[ur, hook]
    &
\end{tikzcd}
\]
where the fibre product is taken along
\[
    F_0:A'_0\to A_0
    \qquad\text{and}\qquad
    p_{V_{\mathbb A}}:V_{\mathbb A}\to A_0.
\]
Including the carrier slices, the condition is
\[
\begin{tikzcd}[column sep=large, row sep=large]
    A'_0\times_{A_0}V_{\mathbb A}
        \arrow[r]
        \arrow[d]
        \arrow[dr, phantom, "\lrcorner", very near start]
    &
    V_{\mathbb A}
        \arrow[d, hook]
    \\
    A'_0\times_{A_0}\mathbf{Beh}_{\mathbb A,\mathbb B,0}
        \arrow[r, "F^\sharp"']
        \arrow[d]
    &
    \mathbf{Beh}_{\mathbb A',\mathbb B,0}
    \\
    Z_{\mathbb A'}
        \arrow[ur, phantom, "\text{over }Z_{\mathbb A'}" description]
    &
\end{tikzcd}
\]
together with factorization through \(V_{\mathbb A'}\).

\begin{definition}[Global models]
\label{def:global-models}
    Let
    \[
        V=(V_{\mathbb A})_{\mathbb A\in\mathsf{Inp}}
    \]
    be a global coequational theory. Define
    \[
        \mathsf{Mod}^{g}(V)_{\mathbb A}
        =
        \{P\in\mathsf{Mly}^{\mathrm{str}}_{\mathbb A,\mathbb B}
          \mid
          \beta_P\text{ factors through }V_{\mathbb A}\}.
    \]
\end{definition}

\subsection{Global closure}
\label{subsec:global-coequational-closure}

\begin{theorem}[Global coequational closure]
\label{thm:global-coequational-closure}
    Let
    \[
        V=(V_{\mathbb A})_{\mathbb A\in\mathsf{Inp}}
    \]
    be a global coequational theory. Then
    \[
        \mathsf{Mod}^{g}(V)
    \]
    is closed under strict homomorphic preimages, regular-epimorphic quotients,
    small coproducts, and input pullback along functors in \(\mathsf{Inp}\).

    If \(V\) is a global \(\mathbb T\)-coequational theory, then
    \[
        \mathsf{Mod}^{g}(V)
    \]
    is also closed under the operations induced by \(\mathbb T\).
\end{theorem}

\begin{proof}
    \leavevmode
    \begin{enumerate}
        \item \textbf{Apply local closure pointwise.}
        Closure under strict homomorphic preimages, regular-epimorphic
        quotients, and small coproducts is
        Theorem~\ref{thm:local-coequational-closure} applied to each
        \(V_{\mathbb A}\). In the structured case, closure under
        \(\mathbb T\)-operations is Theorem~\ref{thm:T-coequational-closure}
        applied pointwise.

        \item \textbf{Check input pullback.}
        Let
        \[
            F:\mathbb A'\to\mathbb A
        \]
        and let
        \[
            P\in\mathsf{Mod}^{g}(V)_{\mathbb A}.
        \]
        Choose a factorization
        \[
            \bar\beta_P:P\to V_{\mathbb A}
        \]
        of \(\beta_P\). Proposition~\ref{prop:input-pullback-mealy}
        gives
        \[
            \beta_{F^*P}
            =
            F^\sharp
            \circ
            \langle p_{F^*P},\,\beta_P\pi_P\rangle.
        \]
        Since \(V\) is global, \(F^\sharp\) sends the pullback of
        \(V_{\mathbb A}\) into \(V_{\mathbb A'}\). Therefore
        \[
            \beta_{F^*P}
        \]
        factors through \(V_{\mathbb A'}\), so
        \[
            F^*P\in\mathsf{Mod}^{g}(V)_{\mathbb A'}.
        \]
        The factorization is
        \[
        \begin{tikzcd}[column sep=large, row sep=large]
            F^*P
                \arrow[r, "{\langle p_{F^*P},\,\bar\beta_P\pi_P\rangle}"]
                \arrow[rrr, bend left=18, "\beta_{F^*P}"]
            &
            A'_0\times_{A_0}V_{\mathbb A}
                \arrow[r, dashed]
            &
            V_{\mathbb A'}
                \arrow[r, hook]
            &
            \mathbf{Beh}_{\mathbb A',\mathbb B,0}.
        \end{tikzcd}
        \]
    \end{enumerate}
\end{proof}

\subsection{Running specializations}
\label{subsec:global-theories-running-specializations}

For Boolean languages, a global coequational theory assigns to each alphabet
\(\Sigma\) a derivative-closed collection
\[
    V_\Sigma\subseteq 2^{\Sigma^*}
\]
such that every monoid homomorphism
\[
    h:\Gamma^*\to\Sigma^*
\]
sends \(V_\Sigma\) into \(V_\Gamma\) by inverse image:
\[
    L\in V_\Sigma
    \quad\Longrightarrow\quad
    h^{-1}L\in V_\Gamma.
\]
If the theories are Boolean, each \(V_\Sigma\) is also closed under the
pointwise Boolean operations.

In the stateless case, globality says that whenever an internal functor
\[
    G:\mathbb A\to\mathbb B
\]
satisfies the coequations at input \(\mathbb A\), then every input
precomposition
\[
    G\circ F:\mathbb A'\to\mathbb B
\]
satisfies the coequations at input \(\mathbb A'\).

\section{The global coequational Birkhoff theorem}
\label{sec:global-coequational-birkhoff}

\subsection{The theorem}
\label{subsec:global-coequational-birkhoff-theorem}

\begin{theorem}[Global coequational Birkhoff theorem]
\label{thm:global-coequational-birkhoff}
    Let
    \[
        \mathcal C=(\mathcal C_{\mathbb A})_{\mathbb A\in\mathsf{Inp}}
    \]
    be a replete global family of full subcategories
    \[
        \mathcal C_{\mathbb A}
        \subseteq
        \mathsf{Mly}^{\mathrm{str}}_{\mathbb A,\mathbb B}.
    \]
    Then \(\mathcal C\) is of the form
    \[
        \mathcal C_{\mathbb A}
        =
        \mathsf{Mod}^{g}(V)_{\mathbb A}
    \]
    for a global coequational theory
    \[
        V=(V_{\mathbb A})_{\mathbb A\in\mathsf{Inp}}
    \]
    if and only if each \(\mathcal C_{\mathbb A}\) is closed under:
    \begin{enumerate}
        \item strict homomorphic preimages;
        \item regular-epimorphic quotients;
        \item small coproducts;
    \end{enumerate}
    and the family \(\mathcal C\) is closed under input pullback along functors
    in \(\mathsf{Inp}\).

    If \(\mathbb B\) is \(\mathbb T\)-structured, then global
    \(\mathbb T\)-coequational theories correspond to such families that are
    also closed under the operations induced by \(\mathbb T\).
\end{theorem}

\begin{proof}
    \leavevmode
    \begin{enumerate}
        \item \textbf{Prove necessity.}
        If
        \[
            \mathcal C=\mathsf{Mod}^{g}(V),
        \]
        then the desired closure properties are
        Theorem~\ref{thm:global-coequational-closure}.

        \item \textbf{Construct the local theories from the model family.}
        Conversely, suppose the displayed closure properties hold. For each
        \(\mathbb A\), define
        \[
            \mathcal S_{\mathcal C,\mathbb A}
            :=
            \{
                S\hookrightarrow
                \mathbf{Beh}^{\mathrm{can}}_{\mathbb A,\mathbb B}
                \mid
                S=\operatorname{Beh}(P)
                \text{ for some }P\in\mathcal C_{\mathbb A}
            \}.
        \]
        Choose a small representative set for
        \(\mathcal S_{\mathcal C,\mathbb A}\), and define
        \[
            V_{\mathcal C,\mathbb A}
            :=
            \bigvee_{S\in\mathcal S_{\mathcal C,\mathbb A}}
            S
            \hookrightarrow
            \mathbf{Beh}^{\mathrm{can}}_{\mathbb A,\mathbb B}.
        \]
        By Theorem~\ref{thm:local-coequational-birkhoff},
        \(V_{\mathcal C,\mathbb A}\) is a local coequational theory and
        \[
            \mathcal C_{\mathbb A}
            =
            \mathsf{Mod}(V_{\mathcal C,\mathbb A}).
        \]

        \item \textbf{Check global reindexing on semantic images.}
        Let
        \[
            F:\mathbb A'\to\mathbb A
        \]
        be an internal functor and let
        \[
            S\in\mathcal S_{\mathcal C,\mathbb A}.
        \]
        Choose
        \[
            P_S\in\mathcal C_{\mathbb A}
        \]
        with
        \[
            S=\operatorname{Beh}(P_S).
        \]
        Closure under input pullback gives
        \[
            F^*P_S\in\mathcal C_{\mathbb A'}.
        \]
        By Lemma~\ref{lem:input-pullback-semantic-images},
        \[
            \operatorname{Im}
            \left(
                A'_0\times_{A_0}S
                \xrightarrow{F^\sharp}
                \mathbf{Beh}_{\mathbb A',\mathbb B,0}
            \right)
            =
            \operatorname{Beh}(F^*P_S).
        \]
        Since
        \[
            F^*P_S\in\mathcal C_{\mathbb A'},
        \]
        we have
        \[
            \operatorname{Beh}(F^*P_S)
            \leq
            V_{\mathcal C,\mathbb A'}.
        \]
        Hence \(F^\sharp\) sends the pullback of each representative semantic
        image
        \[
            S\in\mathcal S_{\mathcal C,\mathbb A}
        \]
        into \(V_{\mathcal C,\mathbb A'}\). The semantic-image step is
\[
\begin{tikzcd}[column sep=large, row sep=large]
    A'_0\times_{A_0}S
        \arrow[rr, "F^\sharp"]
        \arrow[dr, dashed, "\overline{F^\sharp_S}"']
    &&
    \mathbf{Beh}_{\mathbb A',\mathbb B,0}
    \\
    &
    \operatorname{Beh}(F^*P_S)
        \arrow[ur, hook, "m_{F^*P_S}"']
        \arrow[d, hook, "\iota_S"]
    &
    \\
    &
    V_{\mathcal C,\mathbb A'}
        \arrow[uur, hook, "j_{V_{\mathcal C,\mathbb A'}}"']
    &
\end{tikzcd}
\]

        \item \textbf{Descend from the join.}
        Since
        \[
            V_{\mathcal C,\mathbb A}
            =
            \bigvee_{S\in\mathcal S_{\mathcal C,\mathbb A}}S,
        \]
        and pullback preserves joins in a topos slice, the domain
        \[
            A'_0\times_{F_0,A_0,p_{V_{\mathcal C,\mathbb A}}}
            V_{\mathcal C,\mathbb A}
        \]
        is covered by the coproduct
        \[
            \coprod_{S\in\mathcal S_{\mathcal C,\mathbb A}}
            A'_0\times_{A_0}S.
        \]
        On each such piece, \(F^\sharp\) lands in
        \(V_{\mathcal C,\mathbb A'}\). By
        Lemma~\ref{lem:regular-epi-descent-for-subobjects-membership},
        \(F^\sharp\) sends all of \(V_{\mathcal C,\mathbb A}\) into
        \(V_{\mathcal C,\mathbb A'}\). The descent diagram is
        \[
        \begin{tikzcd}[column sep=large, row sep=large]
            \displaystyle
            \coprod_S A'_0\times_{A_0}S
                \arrow[r, two heads]
                \arrow[dr]
                \arrow[ddr, bend right=18, "F^\sharp"']
            &
            A'_0\times_{A_0}V_{\mathcal C,\mathbb A}
                \arrow[d, dashed]
                \arrow[dd, bend left=18, "F^\sharp"]
            \\
            &
            V_{\mathcal C,\mathbb A'}
                \arrow[d, hook]
            \\
            &
            \mathbf{Beh}_{\mathbb A',\mathbb B,0}.
        \end{tikzcd}
        \]
        Thus
        \[
            V_{\mathcal C}
            =
            (V_{\mathcal C,\mathbb A})_{\mathbb A\in\mathsf{Inp}}
        \]
        is a global coequational theory.

        \item \textbf{Handle the structured case.}
        If \(\mathbb B\) is \(\mathbb T\)-structured and the family
        \(\mathcal C\) is closed under the operations induced by
        \(\mathbb T\), then
        Theorem~\ref{thm:T-coequational-birkhoff} applied pointwise shows that
        each \(V_{\mathcal C,\mathbb A}\) is a local
        \(\mathbb T\)-coequational theory. Thus \(V_{\mathcal C}\) is a global
        \(\mathbb T\)-coequational theory.
    \end{enumerate}
\end{proof}

\subsection{Running specializations}
\label{subsec:global-birkhoff-running-specializations}

In the Boolean language specialization, the global theorem recovers the
coequational form of the inverse-image closure condition in Eilenberg-type
correspondences. A family
\[
    (V_\Sigma)_\Sigma
\]
of derivative-closed Boolean language theories is global precisely when
\[
    L\in V_\Sigma
    \quad\Longrightarrow\quad
    h^{-1}L\in V_\Gamma
\]
for every homomorphism
\[
    h:\Gamma^*\to\Sigma^*.
\]
The associated model classes are exactly the recognizers whose generated
residual languages lie in the corresponding \(V_\Sigma\).

In the stateless specialization, the same theorem says that a globally
coequational class of internal functors is determined by its residual slice
functors and is stable under precomposition by input functors. Thus the global
condition is the internal-categorical analogue of inverse morphic preimage.

\section{Summary}
\label{sec:coequations-birkhoff-summary}

This chapter has reformulated residual semantics in terms of coequational
subobjects of the terminal behaviour Mealy morphism.

\[
\begin{array}{c|l}
\text{construction} & \text{role} \\
\hline
\mathbf{Beh}^{\mathrm{can}}_{\mathbb A,\mathbb B}
&
\text{terminal residual semantic Mealy morphism}
\\
V\hookrightarrow \mathbf{Beh}^{\mathrm{can}}
&
\text{local coequational theory}
\\
\beta_P:P\to\mathbf{Beh}^{\mathrm{can}}
&
\text{unique semantic map to the terminal object}
\\
\operatorname{Beh}(P)
&
\text{regular semantic image of a Mealy morphism}
\\
\mathsf{Mod}(V)
&
\text{Mealy morphisms whose semantics land in }V
\\
V_{\mathcal C}
=
\bigvee_{P\in\mathcal C}\operatorname{Beh}(P)
&
\text{coequational theory generated by a model class}
\\
\theta_{\mathbf{Beh}}
&
\text{pointwise structured operation on behaviours}
\\
F^\sharp
&
\text{input reindexing of residual behaviours}
\\
F^*P
&
\text{input pullback of a Mealy morphism}
\end{array}
\]

The local coequational Birkhoff theorem identifies local coequational theories
with replete full classes of strict Mealy morphisms closed under strict
homomorphic preimages, regular-epimorphic quotients, and small coproducts. The
structured version adds closure under pointwise algebraic operations induced by
a \(\mathbb T\)-structured output category. The global version adds stability
under input pullback, equivalently under residual behaviour reindexing.

The next chapter extracts the algebraic side of this coequational semantics.
For a local coequational theory \(V\), the restricted canonical object on \(V\)
determines a residual transition--output action of
\[
    \mathbb A^{\mathrm{op}},
\]
reflecting the target-to-source processing convention. The regular image of
this action is the transition--output quotient associated to \(V\). For a
specification \(K\), applying this construction to the derivative theory
\[
    \operatorname{Der}_0(K)
\]
produces the syntactic transition--output category.